\section{Statistics for the genuine paths}\label{sec:statistics}

The cavity equations require moments of the random overlap $Q$ and
of its centered version $Q-S_v$.  We derive them from two outputs of
the comparison argument.  The mass-averaged budget \eqref{eq:C1}
first gives deterministic averaging windows for the single-path
observable $L_v$.  A useful window must contain enough field increase
for projection while keeping the variation of $S_v$ small.  Balancing
these requirements gives the pointwise and integrated low-mass
estimates below, including the direct-path estimates needed during
the strength comparison.

On size paths, the pointwise scalar-trial excess in \eqref{eq:C2}
has a separate role: it gives uniform coarse accuracy in every fixed
moment.  That argument uses the single-path inequalities and the scalar
gap, and does not use the refined low-mass estimates.  Its power saving
therefore fixes the Taylor order before the final choice of $p_*$.
We then combine the refined low-mass estimates with the pair inverse
to close the high-mass covariance equation.  Throughout, $S_v$ is
the untilted deterministic mean, including when a tilted law is tested.

\subsection{Scales and the pointwise estimates}

Fix the path parameters and the site count $m$.  Unless stated
otherwise, expectations and $L^p$ norms in this section are then under
the genuine untilted positive law.  Strength variations use the direct
path with multipliers $(0,1)$.  The comparison budget \eqref{eq:C1}
allows us to choose a deterministic number $1\le Z\le x^{-C}$, depending
on these fixed data, such that $\mathbf E_\gamma Z^4\lesssim_\circ1$ and,
with $Y=UZ$,
\begin{equation}
\int_{v\asymp r} (S_v-q(v))^2\,dv/r\le C Y^4/r^2\qquad(r\gtrsim U).       \tag{D1}\label{eq:D1}
\end{equation}
For example, let $B$ be the maximum dyadic budget appearing in
\eqref{eq:C1} and take $Z=1+C(B/U^4)^{1/4}$.  This explains both the
fourth parameter moment and the pointwise bound \eqref{eq:D1}.
The same estimate holds on fixed comparable enlargements of the dyads.
Since $(S_v-q(v))^2\le A_v$, the budget indeed controls the mean error.
The number $Z$ is random only when the parameters $\gamma$ are averaged;
it is fixed under every positive-law expectation below.

We first assume $Y=o(1)$, with $U=x\lambda^{-s_*/2}$.  Large $Y$ and
multiplier changes are treated separately.  We will also use a
\textbf{low-region variant}: the path is a size path,
$Y=UZ=o(1)$ satisfies \eqref{eq:D1} at small masses, and refined estimates
are requested only below a sufficiently small fixed mass.  Put
\begin{equation*}
 z_r^0=\frac{x^3}{\sqrt{\lambda_1}\,r^3},\qquad
 g_0=\lambda^{(3s_*-1-\eta)/2},\qquad g_B=g_0/Z^{3/2}.
\end{equation*}
Thus $z_r^0=g_0(U/r)^3$.  These are dimensionless projection scales,
not scalar fields.

In envelope estimates with arbitrarily small power losses, the envelope itself may be chosen for the desired loss: more precisely we use a geometric slack \(b_\epsilon=x^\epsilon\) with \(\epsilon>0\) arbitrarily small, and lose fixed powers of \(b_\epsilon\), as well as subpower factors. Power losses in the resulting applications, involving any fixed moment/expansion orders, can thus all be taken as small as desired. We write \(\lesssim_\diamond\) with this convention. Exponents in such estimates below do not cost an extra power of \(Z\) unless explicitly shown.

\begin{proposition}[Pointwise moments]\label{prop:statistics-pointwise}
Assume the direct-path hypotheses above and $Y=o(1)$, or their low-region
size-path variant.  For every fixed $1\le p<\infty$ the following estimates hold:
\begin{equation}
\begin{array}{ll}
 \||X_v|^2\|_p+\|L_v\|_p+\|Q_{\rm split\ at\ }v\|_p\lesssim_p v+Y;\\
 \|Q_{\rm split\ at\ }v-q(v)\|_p+
 \|2L_v-q(v)\|_p+\||X_v|^2-q(v)\|_p\lesssim_{p,\diamond} e(v)
       & (v\gtrsim Y).
\end{array}                                                        \tag{D2}\label{eq:D2}
\end{equation}
In the low-region variant the second line is only low. \(e\) is deterministic conditional on the parameters, with, on each such dyad \(v\asymp r\),
\begin{equation*}
x^C\le e(v)\lesssim r,\qquad
 \int e(v)^2\,dv/r\lesssim_\diamond m_r^2,\qquad m_r=Y^2/r.
\end{equation*}
On those regions \(|d(v)|+|\dot d(v)|\lesssim_\diamond\lambda r^2 e(v)\), by the corresponding path estimates. At \(v\lesssim Y\) use \(r=e=Y\) in this last assertion and \eqref{eq:D2}.
\end{proposition}

The envelope $e$ depends on the deterministic path data and on the
comparison majorant $Z$, not on a sample from the positive path law.
We first construct it, including the windows needed in the proof.

\subsection{Deterministic widths and averaging windows}
All window lengths in this construction are measured in $z$;
mass integrals are converted by $dv=G_z\,dz$.  On the active interval
of a direct path, set $z=\xi\in[0,z_*]$, where $z_*=\xi_*$.
Extend this coordinate increasingly over the plateau by
\begin{equation*}
 dz=U^2[(1-w+d_+)^{-3/2}+(w-p_{\rm top}+d_-)^{-3/2}]\,dw,
 \qquad w=v/\ell,
\end{equation*}
and write $G_z=dv/dz$.  The function $q(v(z))$ is Lipschitz.  On the
active interval $G$ has bounded derivatives of every fixed order and
is bounded above and below near zero, although it may vanish elsewhere.
On the plateau put
\[
 \mathfrak h(z)=\max\left(
       \frac{U^2}{\sqrt{w-p_{\rm top}+d_-}},
       \frac{U^2}{\sqrt{1-w+d_+}}\right).
\]
The explicit coordinate derivative makes this width Lipschitz and
gives $G_z\mathfrak h^3\asymp U^4$.  At the two ends of the plateau
its width is comparable respectively to $E_b$ and to one.

The extended $z$ interval has bounded length.  Partition it into dyadic
bands $z\asymp r$ near zero and finitely many macroscopic intervals
above, placing a boundary at the active endpoint.  On each band its
length is comparable to $r$, and
$r\asymp\min(z,1)\asymp v$.  Fixed comparable enlargements may be
used for buffers.  We work on bands $r\gtrsim Y$.
On their active parts \eqref{eq:P4} gives a lower field slope, with
the negative ramp slope absorbed as in \eqref{eq:P7}.  On their plateau
parts the formulas for $dz/dw$ and $H_e$ give $dh/dz\gtrsim\lambda_1$;
these masses are macroscopic.  Thus throughout the construction,
\begin{equation}
 dh\ge c\lambda_1r^2\,dz.                                        \tag{D3}\label{eq:D3}
\end{equation}
Write $f(v)=S_v-q(v)$ here; this $f$ is not a pressure.
Monotonicity, \eqref{eq:D1}, and $q(v)\lesssim v$ imply
$S_v\lesssim v+Y$.

The active width has three requirements: enough mass for projection,
enough room for the mean error, and compatibility with the endpoint
ramp.  The next four functions impose these requirements.  Put \(r(z)=\min(z,1)\). Reflect the active \(G_z\) evenly at the endpoints for defining integrals outside, and add an artificial constant \(x^K\), \(K\) arbitrarily large fixed. Use
\begin{equation*}
H_0(z)=\min(b_\epsilon r(z),\,l_0),\qquad
 l_0^2\int_{z-l_0}^{z+l_0}(G_{s,\mathrm{refl}}+x^K)ds
                     =U^6/r(z)^3.
\end{equation*}
Add the cone envelope
\begin{equation*}
H_f(z)=\sup_{s\ {\rm active}}
       [\min(b_\epsilon r(s), |f(v(s))|/Z^2)-A|s-z|]_+
\end{equation*}
with a sufficiently large constant \(A\). Near the upper active endpoint, writing \(d_*=z_*-z,\ E=E_b\), add also the two functions
\begin{equation*}
(E/2-d_*)_+,\qquad
 [\,2\min(E,b_\epsilon)-(d_*-E)_+\,]_+ .
\end{equation*}
Take the maximum of the four functions as \(\mathfrak h(z)\). We only need this width at \(z\gtrsim Y\). Finally set \(e(v(z))=\min(r(z),Z^2\mathfrak h(z))\).

We will repeatedly compare the mass of neighboring intervals in the
$z$ coordinate.  On the active side, if $J$ and $J'$ have comparable
length $l\le Cb_\epsilon$, lie within $O(l)$ of one another, and stay
inside a fixed comparable band, then
\[
 \sup_J G\le C\left(\frac1{|J'|}\int_{J'}G\,dz+x^K\right).
\]
Here and below $K$ may be increased to make additive errors negligible.
To prove this, expand the smooth nonnegative function $G$ to a fixed
order at the scale $l$.  Choose that order so that the Taylor remainder
is at most $x^K$.  On the rescaled bounded intervals, polynomial norm
equivalence bounds the supremum by the average on any comparable
subinterval; nonnegativity permits replacing the polynomial's absolute
average by that of $G$, with the same remainder.  The constants depend
on the fixed Taylor order, not on the interval.  For a reflected
interval meeting an endpoint, restrict to its part on the original
side, whose length is comparable.  This gives the same comparison.

\begin{lemma}[Width construction]\label{lem:statistics-widths}
The preceding construction gives a deterministic envelope with the
size, integral, and defect bounds stated in
Proposition~\ref{prop:statistics-pointwise}.  The random moment bounds
\eqref{eq:D2} will be proved in the next subsection.  The widths and
windows have the following additional properties.
\begin{enumerate}[label=\textup{(\roman*)},leftmargin=*]
\item At comparable points, \(\mathfrak h(z')\lesssim\mathfrak h(z)+|z'-z|\), with comparability of widths when \(|z'-z|\) is a sufficiently small fraction of \(\mathfrak h(z)\). Also \(x^C\le\mathfrak h\lesssim r\), \(e\gtrsim b_\epsilon Y^2/r\), \(|f|\lesssim b_\epsilon^{-1}e\), and \(\int_{\rm band}G_z(\mathfrak h/r)^2 dz/r\lesssim_\diamond(U/r)^4\).
\item Unless \(Z^2\mathfrak h\gtrsim b_\epsilon r\), there are past and future windows of \(z\)-length comparable to \(\mathfrak h\), adjacent or at that distance, each with average \(G_*\) there satisfying
\begin{equation}
\sup G_z\lesssim G_*,\qquad G_*(\mathfrak h/r)^3\gtrsim (U/r)^6 .          \tag{D4}\label{eq:D4}
\end{equation}
They can be taken small fixed fractions of the indicated scales. From such a nontrivial point there is also remaining forward \(z\)-length at least a small fixed macroscopic constant (for taking longer layers up to a macroscopic mass), with \eqref{eq:D3} holding using each successive scale.
\item Partition any full band as above into consecutive tiles of lengths comparable to \(\mathfrak h\), with this width and \(e\) comparable over each tile. Writing \(\mu_I=\int_I G_z dz/r\),
\begin{equation}
\sup_I G_z\lesssim b_\epsilon^{-C}(r\mu_I/|I|+x^K),\quad
 \Delta_I S\lesssim b_\epsilon^{-1}\min(r,Z^2 |I|),\quad
 \sum_I \sup_I(e/r)^2\lesssim Z^2 .                                 \tag{D5}\label{eq:D5}
\end{equation}
The variation assertion holds also on comparable-scale enlargements of these tiles within the path.

\end{enumerate}
\end{lemma}

\Needspace{6\baselineskip}
\begin{proof}
\leavevmode\par
\paragraph{The intrinsic width.}
The map
$l\mapsto l^2\int_{z-l}^{z+l}(G_{s,\mathrm{refl}}+x^K)ds$
is strictly increasing.  Ball inclusion therefore compares its solution
at neighboring centers.  Together with the cap $b_\epsilon r$, this
proves the stated quasi-Lipschitz bound for $H_0$ and comparability
within a small fraction of $H_0$.  Boundedness of $G$ gives
\[
 H_0\gtrsim\min(b_\epsilon r,U^2/r).
\]
The adjacent-interval comparison at this radius also gives
$G_zH_0^3\lesssim U^6/r^3$, up to an arbitrarily small additive error.
Since $G$ is bounded, the interpolation
\[
 G_zH_0^2=(G_zH_0^3)^{2/3}G_z^{1/3}\lesssim U^4/r^2
\]
proves the squared width budget after integration over a band of
$z$-length $O(r)$ and division by $r^3$.

\paragraph{The deviation cone.}
For each source point $s$, the function
$[\min(b_\epsilon r(s),|f(v(s))|/Z^2)-A|s-z|]_+$
is $A$-Lipschitz.  Their supremum is therefore Lipschitz too.  At its
source, a height of order $l$ implies $|f|\gtrsim Z^2l$.
If $f$ is positive, monotonicity of $S$ and Lipschitz continuity of
$q(v(z))$ preserve half this deviation on a forward interval of
length $cl$; for a negative deviation use the backward interval.
Thus each source of height $l$ supplies a one-sided interval $J_s$ with
\[
 Z^4l^2\int_{J_s}G_z\,dz
       \lesssim \int_{J_s}f(v(z))^2G_z\,dz.
\]
Its cone is supported within $O(l)$ of $s$.  The adjacent-interval
comparison bounds the mass of that support by the mass of $J_s$, with
negligible additive error.  At each dyadic height $l\ge x^K$, choose
separated source points to cover the cone superlevel set.  The
corresponding intervals have bounded overlap at that height.  Summing
the preceding inequality, and then the logarithmically many heights,
gives
\[
 \int_{\rm band}G_zH_f(z)^2\,dz
 \lesssim_\diamond Z^{-4}\int_{\rm enlarged\ band}f(v)^2\,dv
 \lesssim_\diamond U^4/r.
\]
Sources of height below $x^K$ contribute negligibly.  The lower boundary
causes no difficulty because $l\le b_\epsilon r$.  If a positive
source's forward interval reaches the upper active endpoint, its
deviation persists through the plateau.  The macroscopic version of
\eqref{eq:D1} then gives $Z^4l^2\lesssim Y^4$, or $l\lesssim U^2$;
these sources already obey the same macroscopic squared width budget.

\paragraph{The plateau and its active edge.}
The reference $q$ is constant on the plateau.  A positive
deviation of $S-q$ persists towards its upper end; a negative deviation
persists towards its lower end.  Integrate \eqref{eq:D1} over the
corresponding endpoint strip, whose mass is comparable to the distance
to that endpoint.  Away from the capped strips, the result is
$|f|\lesssim Z^2\mathfrak h$.  Inside the upper $d_+$ strip,
$\mathfrak h\asymp1$ and the bound follows from boundedness of $S,q$.
For a negative deviation inside the lower $d_-$ strip, the backward
interval must instead use
the active side.  The last active $z$-interval of length $E=E_b$ has
mass comparable to $d_-=U^4/E^2$.  Either $|f|=O(Z^2E)$ already, or
the Lipschitz variation of $q$ over that interval is too small to
remove half the negative deviation.  In the latter case \eqref{eq:D1}
gives
\[
 |f|^2\frac{U^4}{E^2}\lesssim Y^4,
 \quad\hbox{hence}\quad |f|\lesssim Z^2E,
\]
which is the required lower-edge bound.

The explicit plateau formula also gives
$\int_{\rm plateau}\mathfrak h^2\,dv\lesssim_\circ U^4$ by
integrating $U^4/(d+d_\pm)$ from each endpoint.  The first extra active
function has height $O(E)$ on a region of mass $O(U^4/E^2)$, so its
squared cost is $O(U^4)$.  The second has height at most
$2\min(E,b_\epsilon)$; its part inside the last $E$ costs the same.
For its additional layer of length $2\min(E,b_\epsilon)$, compare
adjacent active intervals at that length.  Their mass is controlled by
the preceding last-$E$ region, giving again $O_\diamond(U^4)$.
These two functions also majorize the active ramp defect and its
parameter derivative divided by $\lambda$, with at most a
$b_\epsilon^{-1}$ loss.  The bounds in \eqref{eq:P6}--\eqref{eq:P8}
then give the claimed defect estimates.  Below $Y$, those estimates
use only the low and mask bounds from \eqref{eq:P6}, \eqref{eq:P8}.

\paragraph{Matching the two coordinates at the join.}
Suppose first that $E\lesssim b_\epsilon$, and let $d_*=z_*-z$ on
the active side.  We claim, for $d_*\le Cb_\epsilon$, that
\[
 |f(v(z))|/Z^2\lesssim E+d_*.
\]
For $d_*\lesssim E$, use the last active interval of length $E$ or
the adjacent plateau, in the direction in which the deviation
persists.  The preceding mass calculation proves the claim.
For $d_*\gtrsim E$, an interval of length comparable to $d_*$ fits
on the active side in either direction.  Unless the claimed bound
already holds, the deviation persists on that interval.  Adjacent
interval comparison bounds below its mass by that of the last active
length $E$, up to negligible error.  Applying \eqref{eq:D1} as above
would then give $|f|\lesssim Z^2E$, proving the claim in this case too.

With $A$ sufficiently large, sources of the cone obeying this linear
bound imply $H_f(z)\lesssim E+d_*$.  The same bound holds for $H_0$:
a ball of radius $C(E+d_*)$ contains the last active length $E$, and
its squared radius times that mass is at least a constant times $U^4$,
which exceeds the intrinsic target $U^6/r^3$ on this macroscopic band.
When $E>b_\epsilon$, the corresponding upper bounds follow from the
caps in $H_0,H_f$.  The explicit endpoint functions provide the lower
bound near the join.  Thus at distances much smaller than $E$, both
the active and plateau widths are comparable to $E$.  Their separate
quasi-Lipschitz bounds consequently remain valid across the join.

\paragraph{Windows and tiles.}
Consider first an active point with $d_*\ge E$ for which
$Z^2\mathfrak h<b_\epsilon r$ up to the fixed threshold in (ii).
The intrinsic cap is then inactive.  Its defining equation and
$\mathfrak h\ge H_0$ give
\[
 \mathfrak h^2\int_{z-C\mathfrak h}^{z+C\mathfrak h}G_s\,ds
       \gtrsim U^6/r^3,
\]
with negligible additive error.  Choose a past interval and a future
interval, each a small fixed fraction of $\mathfrak h$ and within
$O(\mathfrak h)$ of $z$.  Since
$\mathfrak h\lesssim E+d_*\lesssim d_*$, sufficiently small fixed
fractions fit wholly on the active side.  Near zero the cap
$\mathfrak h\lesssim b_\epsilon r$ also leaves room for the past
window.  Adjacent-interval comparison gives, for each chosen window,
$\sup G\lesssim G_*$ and
$G_*\mathfrak h^3\gtrsim U^6/r^3$, which is \eqref{eq:D4}.

For active points with $0\le d_*\le E$, the endpoint functions
and the nontrivial accuracy requirement force $E\lesssim b_\epsilon$
and $\mathfrak h\asymp E$.  The plateau width is likewise comparable
to $E$ in a sufficiently small fixed $E$-neighborhood of the join.
The mass of the last active interval of length $E$ is comparable to
$U^4/E^2$, so comparable active windows there have average density
$G_*\asymp U^4/E^3$.  The adjoining plateau windows in that small
neighborhood have the same bound.  Hence
$G_*\mathfrak h^3\asymp U^4\gtrsim U^6$, as required on this
macroscopic band.  Outside that neighborhood on the plateau, the
explicit width formula allows small fixed fractional windows wholly
within the plateau and makes $G$ comparable on each of them.  The
identity $G\mathfrak h^3\asymp U^4$ gives the same conclusion.  Finally the terminal $d_+$ strip has
macroscopic $z$-length and width comparable to one.  A point requiring
nontrivial accuracy must precede that strip, leaving the fixed forward
length needed for longer projection layers.

Tiling can for instance use equal bounded small steps in the integral of \(dz/\mathfrak h\) on each band (resizing the steps by bounded factors to fit). Width comparability holds by the local property. On an active tile longer than \(b_\epsilon\), subdivide into \(b_\epsilon\)-size pieces for \eqref{eq:D5}; otherwise use the adjacent-interval bound directly. For a tile $I$ of length comparable to $\mathfrak h$, use
$S=q+f$, the Lipschitz bound on $q(v(z))$, and
$|f|\lesssim b_\epsilon^{-1}e$ on its enlargement.  This gives
\[
 \Delta_I S\lesssim b_\epsilon^{-1}\min(r,Z^2|I|).
\]
Finally, with $a_I=\sup_I e/r$, width comparability and
$e=\min(r,Z^2\mathfrak h)$ imply
\[
 a_I^2\lesssim \frac{Z^2|I|}{r},\qquad
 \sum_I a_I^2\lesssim Z^2,
\]
since the disjoint tiles fill an interval of length $O(r)$.
The same argument applies to fixed comparable enlargements.
Combining the squared budgets for the four defining widths also gives
\[
 \int e(v)^2\,dv/r
 \le Z^4\int\mathfrak h(z)^2G_z\,dz/r
 \lesssim_\diamond Y^4/r^2=m_r^2.
\]
The lower intrinsic width yields $e\gtrsim b_\epsilon Y^2/r$ for
$r\gtrsim Y$, including when either cap is active.

\paragraph{The low-region construction.}
Choose a small fixed working interval and a larger fixed buffer
interval on which the size-path inverse is regular.  Use $z=v$ there,
so $G=1$ and \eqref{eq:D3} holds.  Define
\[
 \mathfrak h=\max\{\min(b_\epsilon r,Y^2/r),H_f\},
 \qquad e=\min(r,\mathfrak h),
\]
where the cone now uses $|f|$ in place of $|f|/Z^2$.
The preceding cone argument gives squared width budget $(Y/r)^4$
on the working interval.  The intrinsic lower width is now
$\min(b_\epsilon r,Y^2/r)$, which is at least the direct-path lower
width because $Y\ge U$.  Consequently the window bound
\eqref{eq:D4} remains valid.  The pointwise and variation bounds use
$e=\min(r,\mathfrak h)$, and the tile sum is now
$\sum_I\sup_I(e/r)^2\lesssim1$; these are precisely the later bounds
in which $Z_h=1$.

If more room is needed for buffers, start a fixed Lipschitz ramp only
beyond the working interval, increasing $\mathfrak h$ to a macroscopic
value before the end of the buffer interval.  No integrated width
estimate is used on that auxiliary ramp.  Longer layers stop after a
fixed small macroscopic length within the regular region.  The defect
bounds are the low size-path bounds of \eqref{eq:P8}.

\paragraph{Measurability of the construction.}
All windows can be chosen measurably and deterministically conditional
on the path parameters.  For example, fix the fractional lengths and
casewise windows constructed above, using the prescribed comparable
band enlargements for buffers, and tile by equal increments of
$\int dz/\mathfrak h(z)$.  The cone
supremum may be formed over a countable dense set together with the
specified one-sided values of the monotone function $S$.  Its value is
unchanged, since $q$ is continuous and the one-sided limits of $S$ exist.
The implicit intrinsic radius is the monotone solution of its displayed
integral equation.  These constructions are measurable functions of
the deterministic path data.  In particular, none of the windows is
selected after observing a spin sample.
\end{proof}

\Needspace{8\baselineskip}
\subsection{Proof of the pointwise moments}
\begin{proof}[Proof of Proposition~\ref{prop:statistics-pointwise}]
All estimates in this proof use the untilted positive law with fixed
path data. Lemma~\ref{global:residual-blocks} supplies the conditional
block decomposition \eqref{eq:global-block-decomposition} and its
integrated-increment bounds. At each deterministic fork cut,
partition the residual on each of two conditionally independent
continuations into future martingale increments, including a last residual ending at the terminal
Gibbs spin. A block pair contributes its squared inner product,
weighted by the largest mass in the Hessian block. We may bound that
weight by the largest mass in either block and project the block with
the later starting point.

For the opposite increment write $A=|A|\widehat A$, setting
$\widehat A=0$ when $A=0$. Conditional on the common prefix, its
whole branch is exterior data for the projected continuation. Applying
\eqref{eq:A3} to $\widehat A$ therefore preserves the projected
branch's normalized future kernel. The factor $|A|^2$ remains in the
subsequent integral. Fix a permitted loss $\varepsilon>0$. At each
fixed fork, declare an exception when the squared projection exceeds
$x^{-\varepsilon}$ times its squared scale in \eqref{eq:A3}.
Choose the projection moment sufficiently high after $\varepsilon$
and the required negligible exponent $K$ are fixed. Markov's
inequality then makes this exception have probability $O(x^K)$.
The factor $x^{-\varepsilon}$ is included in the permitted losses
below. Bounded spins,
integration over the fork coordinate, and conditional Jensen make its
contribution negligible in the integral norm. This does not require an
event simultaneous over all real forks, or a higher-moment bound for
$|A|$ proportional to the bootstrap constant. These operations can
first be performed with inserted deterministic cuts and then passed to
the fixed-size path limit.

\paragraph{Raw-moment bootstrap.}
Fix a large even exponent $p$ and set
\[
 B_0=\max\left(1,\sup_v\frac{\||X_v|^2\|_p}{v+Y}\right).
\]
This number is finite at each fixed size because $|X_v|\le1$ and
$Y>0$. At a small scale $r\gtrsim Y$ in the regular inverse region,
let $J$ be a fork window of $z$-length comparable to $r$. We claim
\begin{equation}
 \left\|\int_J W_z\,dz/r\right\|_{p/2}
 \lesssim_{p,\circ}(z_r^0)^2r^3B_0.
 \tag{D6}\label{eq:D6}
\end{equation}

Use an initial block of arbitrarily small inverse-polynomial length,
then blocks starting at doubled distances from the fork, and finally
a terminal residual. Stop the doubling at a fixed small macroscopic
distance, with a full preceding interval still inside the regular
region. Consider a block pair other than two initial blocks, and let
$lr$ be its later starting distance. If $l\lesssim1$, the projection
gap has mass at least $cr$ and field length at least
$c\lambda_1r^3l$. Its squared projection coefficient, before the mass
weight, is consequently
$O_\circ(x^6/(\lambda_1r^5l))$.

The earlier squared anchor has the integral bound
\[
 \left\|\int_J|A_z|^2\,dz/r\right\|_{p/2}
 \le C_p lB_0r.
\]
Indeed, \eqref{eq:global-grid} integrates translated grids of disjoint
increments, and the raw bootstrap bounds the terminal squared norm on
the enlarged window by $CB_0r$. If that anchor belongs to the
independently resampled continuation, apply
\eqref{eq:global-dual-doob} to its nonnegative squared increment to
average it back to the fork. Thus the contribution from this block
pair is at most
\[
 \underbrace{Cr}_{\text{mass weight}}\,
 \underbrace{C_{p,\circ}\frac{x^6}{\lambda_1r^5l}}
             _{\text{squared projection coefficient}}\,
 \underbrace{C_p lB_0r}_{\text{integrated squared anchor}}
 \lesssim_{p,\circ}(z_r^0)^2r^3B_0.
\]
The two initial blocks instead use the sliding estimate without
projection. Choose their length so small that the resulting term is
negligible after the polynomial prefactors.

For a later starting distance comparable to $r'\gtrsim r$, use a
full gap at masses comparable to $r'$, with field length at least
$c\lambda_1(r')^3$. This includes a macroscopic $r'$ for the
terminal block. The mass weight, squared projection coefficient, and
raw squared anchor now cost respectively
\[
 Cr',\qquad
 C_{p,\circ}\frac{x^6}{\lambda_1(r')^5},\qquad
 C_pB_0r'.
\]
Their product is
$O_{p,\circ}((z_{r'}^0)^2(r')^3B_0)$, no larger than the
right side of \eqref{eq:D6}. The normalized fork window has bounded
length. Summing the logarithmically many block pairs in
\eqref{eq:global-block-decomposition} proves \eqref{eq:D6}. The
sliding and conditional integration bounds are uniform in the inserted
grid sizes, so this argument uses no size-uniform path approximation
rate and no derivative of $S$.

Let $L_J=|J|^{-1}\int_JL_u\,du$ for a future mass window
$J$ of width comparable to $r$.  The future bracket identity in
Lemma~\ref{global:buffer-brackets} gives
$|X_v|^2\le2\mathbf E_vL_J$, while $\mathbf EL_J=O(r)$ by the
deterministic mean bound.  Equations \eqref{eq:A1} and \eqref{eq:D6}
give
\[
 \|L_J-\mathbf EL_J\|_p^2
 \lesssim r^{-1}\left\|\int_JW_u\,du/r\right\|_{p/2}
 \lesssim_{p,\circ}(z_r^0)^2r^2B_0.
\]
Conditional Jensen therefore yields
$\||X_v|^2\|_p\lesssim_pr(1+z_r^0\sqrt{B_0})$.
For $v\lesssim Y$ use a future window at scale $Y$; at macroscopic
masses use $|X_v|\le1$.  Taking the supremum in the definition of $B_0$
now gives $B_0\le C_p(1+g_0\sqrt{B_0})$, including the harmless
subpower factors.  The fixed power saving in $g_0$ absorbs these
factors, and this quadratic inequality bounds $B_0$ uniformly.

For \(L,Q\), project endpoints into a future interval at comparable scale \(v+Y\) and then telescope both branches in future doubled layers with terminal block, as above. Each tail product uses \eqref{eq:A3} on the later layer (with comparable preceding clock fitting strictly after the fork); the opposite anchor has norm \(O_p(\sqrt{r'})\) at worst at scale \(r'\), so these tails cost \(O_{p,\circ}(z_{r'}^0 r')\). For \(L\) the anchor \(X_v\) already precedes the gap. This gives the first line of \eqref{eq:D2}. It suffices here too to telescope at small scales.

\paragraph{Refined-moment bootstrap.}
For this argument write $Z_h=Z$ on direct paths and $Z_h=1$ in the
low-region size-path variant. Let $B_1\ge1$ be the least constant for
which
\begin{equation*}
 \||X_{z'}-X_z|^2\|_p
 \le B_1Z_h^2(|z-z'|+\mathfrak h(z))
\end{equation*}
holds for comparable nearby cuts in the bands $r\gtrsim Y$ and their
buffers. The supremum defining $B_1$ is finite at fixed size because
$\mathfrak h\ge x^C$; take it simultaneously over the bands to avoid
artificial interior boundaries. At $z\asymp Y$ the raw bound already
suffices with a $b_\epsilon^{-1}$ loss, as it does on the large-width
auxiliary ramp. Cuts whose separation is not a small fraction of their
scale can likewise use the raw bound.

Let $J$ be one of the windows in \eqref{eq:D4} near $z_0$, and put
$h_0=\mathfrak h(z_0)/r$, so $|J|\asymp rh_0$. We claim
\[
 \left\|\int_J W_z\,dz/r\right\|_{p/2}
 \lesssim_{p,\diamond}(z_r^0)^2r^3h_0B_1Z_h^2.
\]
For a block pair whose later starting distance is $lr$ with
$l\ll1$, the same mass and projection factors as above apply. The
integrated earlier anchor satisfies, in both ranges of $l$,
\[
 \left\|\int_J|A_z|^2\,dz/r\right\|_{p/2}
 \le C_p lB_1Z_h^2rh_0.
\]
When $l\ge h_0$, the bootstrap and quasi-Lipschitz width bound give
$\||A_z|^2\|_{p/2}\le CB_1Z_h^2lr$ pointwise; integration contributes
$|J|/r\asymp h_0$. When $l<h_0$, apply
\eqref{eq:global-grid} after subtracting the initial martingale value
on an enlargement of $J$. Its terminal squared increment is bounded
by $CB_1Z_h^2rh_0$, and the sliding estimate contributes $l$.
The conditional integration inequality handles a resampled anchor in
exactly the same way as in the raw argument. Hence each such pair costs
\[
 \underbrace{Cr}_{\text{mass weight}}\,
 \underbrace{C_{p,\diamond}\frac{x^6}{\lambda_1r^5l}}
             _{\text{squared projection coefficient}}\,
 \underbrace{C_p lB_1Z_h^2rh_0}_{\text{integrated squared anchor}}
 \lesssim_{p,\diamond}(z_r^0)^2r^3h_0B_1Z_h^2.
\]
The pair of initial blocks again has negligible cost. Longer layers
use the already closed raw bound. Their product at scale $r'\gtrsim r$
is $O_{p,\diamond}((z_{r'}^0)^2(r')^3)$ pointwise, and the fork
window contributes $h_0$. Since $B_1Z_h^2\ge1$ and
$(z_{r'}^0)^2(r')^3\le(z_r^0)^2r^3$, these layers obey the same
bound. The forward length in Lemma~\ref{lem:statistics-widths} supplies
their full projection gaps, including the terminal one. Summing the
blocks proves the claim.

The mass of this window is comparable to $G_*rh_0$, and
$\sup G\lesssim G_*$.  Apply \eqref{eq:A1} with the constant weight
$1/(G_*rh_0)$ on the window.  The squared centered norm of its mass
average is at most
\begin{equation*}
 C_{p,\diamond}
 \frac{G_*r\,(z_r^0)^2r^3h_0B_1Z_h^2}{(G_*rh_0)^2}
 = C_{p,\diamond}\frac{(z_r^0)^2r^2B_1Z_h^2}{G_*h_0}
 \lesssim_{p,\diamond}r^2B_1Z_h^2h_0^2.
\end{equation*}
In the last step we used $z_r^0=g_0(U/r)^3$ and
$G_*h_0^3\gtrsim(U/r)^6$ from \eqref{eq:D4}.
We now apply the bracket identities of
Lemma~\ref{global:buffer-brackets}. Earlier and later averages bracket
$|X_z|^2$. To see also
how they control increments, let $s\le z\le z'$ and use
\[
 2\mathbf E_zL_s=|X_z|^2-|X_z-X_s|^2.
\]
Average $s$ over one common past buffer before $z$.  A future buffer
after each of $z,z'$ bounds its squared norm from above; the common
past identity then bounds the averages of $|X_z-X_s|^2$ and
$|X_{z'}-X_s|^2$.  For every such $s$,
\[
 |X_{z'}-X_z|^2
 \le2|X_{z'}-X_s|^2+2|X_z-X_s|^2.
\]
The means of these buffer differences cost at most
$Cb_\epsilon^{-1}Z_h^2(|z'-z|+\mathfrak h(z))$ by the deterministic
variation bound.  Their centered parts have the bound just proved,
with conditional Jensen where needed.  Consequently
\[
 \||X_{z'}-X_z|^2\|_p
 \le C_{p,\diamond}(1+\sqrt{B_1})
       Z_h^2(|z'-z|+\mathfrak h(z)).
\]
If either cut has $Z_h^2\mathfrak h\gtrsim b_\epsilon r$, the raw bound
already gives this inequality, using the quasi-Lipschitz width bound
to compare the two cuts.  The lower band and the auxiliary ramp are
handled in the same way.  The supremum defining $B_1$ therefore
satisfies $B_1\le C_{p,\diamond}(1+\sqrt{B_1})$, so it is bounded with
the permitted loss.  Applying the same two-sided brackets at a single
cut gives $\||X_v|^2-q(v)\|_p\lesssim_{p,\diamond}e(v)$.

Finally project the endpoints of an overlap, or the terminal
endpoint in $L$, to a cut a distance comparable to $\mathfrak h$ later.
Polarization about the common ancestor expresses the projected inner
product through the squared norms and squared differences just bounded.
The remaining tails, estimated by \eqref{eq:A3} and longer layers,
cost at most $O_{p,\diamond}(z_r^0r/\sqrt{h_0})$.  The lower width
$h_0\gtrsim b_\epsilon(U/r)^2$ gives
\[
 \frac{z_r^0r}{\sqrt{h_0}}
 =rh_0\frac{g_0(U/r)^3}{h_0^{3/2}}
 \lesssim b_\epsilon^{-3/2}g_0rh_0
 \lesssim_\diamond e(v)
\]
at every nontrivial cut.  This completes \eqref{eq:D2}.
\end{proof}

\subsection{Integrated variances and centered polynomials}

The pointwise envelope locates the regions of small overlap error.
For the pressure expansion we also need a gain after integration in
mass, including on regions where that envelope is relatively large.
We first use the source inequality and the deterministic variation of
$S$ on each tile to bound the integral of $W$.  Mass averaging then
gives integrated overlap variances, and the polynomial single-path
inequality gives centered integral bounds for the later allocation
argument.

Continue with \eqref{eq:D1}--\eqref{eq:D6}; band integrals without limits below are over \(v\asymp r\), \(r\gtrsim Y\). Put \(Z_h=Z\) for the direct path, \(Z_h=1\) in the low-region variant. The following estimates concern active overlaps (unmasked).
\begin{lemma}[Integrated variance estimates]\label{lem:statistics-integrated}
For every fixed integer $k\ge0$, the following bounds hold. All quantities on the right are untilted; in particular this convention applies to \(S,D\):
\begin{equation}
\begin{split}
 \mathbf E\int W\,e(v)^{2k}\,dv/r
    &\lesssim_\diamond (z_r^0)^2 Z_h^2 r^3
         \begin{cases}1& k=0,\\ m_r^2 r^{2k-2}& k\ge1,\end{cases}\\
 \int \mathbf E[X^{\mathsf T} C_v X+\operatorname{tr}C_v^2]\,dv/r
    &\lesssim_\diamond (z_r^0)^2 Z_h^2 r^2,\qquad
 \int D_v^2\,dv/r\lesssim_\diamond z_r^0 Z_h r^2 .
\end{split}
\tag{D7}\label{eq:D7}
\end{equation}
\end{lemma}
\begin{proof}
On a tile $I$ in a band $v\asymp r$, the inequalities
$v\mathbf EW\le\mathbf E\operatorname{tr}H^2$ and
$dS\ge m\mathbf E\operatorname{tr}H^2\,dh$ from \eqref{eq:A1},
together with \eqref{eq:D3}, imply
\[
 \int_I\mathbf EW_z\,dz/r
 \lesssim \frac{x^6}{\lambda_1r^4}\Delta_I S
 \lesssim_\diamond (z_r^0)^2r^3Z_h^2\frac{|I|}{r}.
\]
Here $Z_h=Z$ for the direct construction and $Z_h=1$ for the low-region
construction, exactly as in their respective variation bounds.
Multiply by the supremum of $G$ on $I$ and use \eqref{eq:D5} to obtain
\begin{equation*}
 \int_I\mathbf EW_zG_z\,dz/r
 \lesssim_\diamond (z_r^0)^2r^3Z_h^2(\mu_I+x^K).
\end{equation*}
There are only polynomially many tiles because their widths have a
polynomial lower bound; increasing $K$ absorbs all additive errors.
Summing gives the first line of \eqref{eq:D7} when $k=0$.  When
$k\ge1$, $e$ is comparable on each tile and $e\lesssim r$, so
\[
 \sum_I\sup_I e^{2k}\,\mu_I
 \lesssim r^{2k-2}\int e^2\,dv/r
 \lesssim_\diamond r^{2k-2}m_r^2.
\]
This proves the weighted version.  Since $H_v\succeq vC_v$,
$\mathbf E\operatorname{tr}C_v^2\le\mathbf EW_v/v$; the same bound
also proves the trace part of the second line.

\paragraph{The mixed covariance term.}
Telescope $X_z$ backwards over dyadic distances $lr$, down
to an arbitrarily small inverse-polynomial length, and retain an oldest
anchor a fixed small fraction of $r$ earlier.  For each nonterminal
increment, the gap after its end allows \eqref{eq:A3} against
$y-X_z$.  Its squared projection coefficient is at most
$C_\circ(z_r^0)^2r/l$, multiplied by the increase of $S$ over
$[z-2lr,z-lr]$.

We need the integral of that increase against $G_z\,dz/r$.
If $lr\gtrsim\mathfrak h(z)$, the deterministic variation estimate
bounds it pointwise by $C_\diamond Z_h^2lr$.  If $lr$ is shorter
than a tile $I$, Fubini for the positive measure $dS$ gives
\[
 \int_I[S(z-lr)-S(z-2lr)]\,dz
       \le lr\,\Delta_{I^+}S,
\]
where $I^+$ is a fixed enlargement.  Multiply by $\sup_I G/r$,
use $\Delta_{I^+}S\lesssim_\diamond Z_h^2|I|$, and sum the tile
mass bounds from \eqref{eq:D5}.  In both cases the weighted integral
is $O_\diamond(rZ_h^2l)$.  Multiplication by the projection
coefficient therefore costs $O_\diamond((z_r^0)^2r^2Z_h^2)$.
The last tiny increment is treated without projection by the same
sliding estimate, choosing its length to dominate all polynomial
prefactors.  The oldest anchor uses a comparable field gap from
\eqref{eq:D3} and its raw moment from \eqref{eq:D2}.
Only logarithmically many layers occur.  Summing proves the claimed
mixed covariance bound.

\paragraph{Mass averaging.}
Let $h=rz_r^0Z_h$ be a width in the mass coordinate.  If $h$
is comparable to or larger than $r$, the raw bound $D_v=O_p(r)$ already
proves the desired estimate.  Otherwise bracket $|X_v|^2$ by mass
averages of $L$ on its two sides.  The points within distance $O(h)$
of a band edge can be discarded at cost $O(r^2h/r)$ using
\eqref{eq:D2}.  For the remaining points, monotonicity and
$\Delta S\lesssim r$ give the same cost for the integrated squared
mean error.  More explicitly, a sliding increment of $S$ is at most
$Cr$, and its integral in $dv/r$ is $O(h)$; its squared integral is
therefore $O(rh)$.  By \eqref{eq:A1} and Fubini, the centered averages
contribute at most $C h^{-1}\int\mathbf EW\,dv/r$.  Thus
\[
 \int\operatorname{Var}(|X_v|^2)\,dv/r
 \lesssim_\diamond r h+
       \frac{(z_r^0)^2Z_h^2r^3}{h}
 \lesssim_\diamond r^2z_r^0Z_h.
\]
The first two parts of the lemma control the other terms in the exact
conditional-covariance identity
\[
 D_v^2=\operatorname{Var}(|X_v|^2)
          +2\mathbf EX_v^{\mathsf T}C_vX_v
          +\mathbf E\operatorname{tr}C_v^2.
\]
This proves the last assertion of \eqref{eq:D7}.

\end{proof}

\begin{lemma}[Large-envelope localization]\label{lem:statistics-localized}
There is also a localized bound for the direct path. For any sufficiently small fixed \(\delta'>0\), on all the bands in \eqref{eq:D7} together,
\begin{equation}
\int_{\{e(v)\ge x^{\delta'}\}} D_v^2\,dv
       \lesssim_\diamond x^{-C_0\delta'} z_1^0 Y^2 Z^2                 \tag{D8}\label{eq:D8}
\end{equation}
with an absolute finite loss exponent \(C_0\) independent of \(\delta'\). \end{lemma}
\begin{proof}
Retain every tile meeting $\{e\ge x^{\delta'}\}$.  Such a tile
has $r\gtrsim x^{\delta'}$, and width comparability makes $e$ at least
a constant times $x^{\delta'}$ throughout it.  Write
$a_I=\sup_I e/r$ and recall that its mass is $r\mu_I$.

Bracket within a retained tile using mass windows of width $h$.
The pointwise error bound in \eqref{eq:D2} gives squared accuracy
$O_\diamond(r^2a_I^2)$ near the tile edges.  Since their mass is
$O(h)$, this costs $O_\diamond(ra_I^2h)$ in $dv/r$.
Monotonicity and $\Delta_IS\lesssim_\diamond ra_I$ give that same
bound for the interior mean errors.  The tile estimate for $W$ in the
preceding proof gives centered cost
$O_\diamond((z_r^0)^2r^3Z^2\mu_I/h)$.  Balance these two terms by taking
\[
 h=\frac{rz_r^0Z\sqrt{\mu_I}}{a_I},\qquad
 ra_I^2h+\frac{(z_r^0)^2r^3Z^2\mu_I}{h}
       \lesssim r^2z_r^0Z a_I\sqrt{\mu_I}.
\]
If $h$ does not fit in the tile, use the raw squared error on the whole
tile instead.  Its cost $r^2a_I^2\mu_I$ is no larger than the displayed
bound when $h\gtrsim r\mu_I$.  This also handles overlapping buffers.
Tiles of negligible mass may first be discarded so the additive
$x^K$ errors in the $W$ bound do not enter the denominator.

The trace contribution is bounded by the tile $W$ estimate divided
by $r$.  For the mixed term use Cauchy--Schwarz and the raw fourth
moment of $X$ on the tile.  These two contributions together cost at
most $O_\diamond(r^2z_r^0Z\mu_I)$.  To sum them and the variance
bound, the squared envelope budget gives on the retained tiles
\[
 \sum_I\mu_I
 \lesssim_\diamond x^{-2\delta'}\frac{Y^4}{r^2}
 \lesssim_\diamond x^{-C\delta'}Y^4,
 \qquad
 \sum_I a_I\sqrt{\mu_I}
 \le\Big(\sum_Ia_I^2\Big)^{1/2}
       \Big(\sum_I\mu_I\Big)^{1/2}
 \lesssim_\diamond Zx^{-C\delta'}Y^2.
\]
The second inequality uses \eqref{eq:D5}.  Converting $dv/r$ back to
$dv$ multiplies a band bound by $r$, and $r^3z_r^0=z_1^0$.
Since $Y=o(1)$ and $Z\ge1$, the sum of the mixed and trace terms is
also bounded by $x^{-C\delta'}z_1^0Y^2Z^2$.  Summing the logarithmically
many retained bands gives \eqref{eq:D8}, with a fixed exponent $C_0$
independent of $\delta'$.

\end{proof}

\begin{lemma}[The lower band]\label{lem:statistics-lower-band}
At \(v\le C Y\) we will discard \(v<x\) or treat it by raw bounds. For dyads of size \(w\lesssim Y\) restricted to \(v\ge x\), put
\(\epsilon_w=x^3\sqrt Y/(\sqrt{\lambda_1}\max(w,U)^{3/2})\), \(E_0=g_B Y^2=x^3\sqrt Y/(\sqrt{\lambda_1}U^{3/2})\). Then
\begin{equation}
\int_{v\asymp w,\ v\ge x}\mathbf EW\,dv\lesssim\epsilon_w^2,\qquad
 \int_{v\asymp w,\ v\ge x} D_v^2\,dv\lesssim_\circ Y\epsilon_w.           \tag{D9}\label{eq:D9}
\end{equation}
\end{lemma}
\begin{proof}
On these dyads,
$v\,dh/dv\gtrsim\lambda_1\max(v,U)^3$ and $S\le CY$.
The source inequality in \eqref{eq:A1}, integrated against the total
increase of $S$, therefore gives
\[
 \int\mathbf EW\,dv
 \lesssim \frac{x^6Y}{\lambda_1\max(w,U)^3}=\epsilon_w^2.
\]
If $w<\epsilon_w/Y$, the raw bound gives directly
$\int D_v^2dv\lesssim Y^2w\le Y\epsilon_w$.  In the other case use
mass slides of width $h=\epsilon_w/Y$.  The mean and boundary errors
cost $O(Y^2h)$, while the centered error costs $O(\epsilon_w^2/h)$;
both equal $O(Y\epsilon_w)$.  Also
$\int\mathbf E\operatorname{tr}C_v^2dv\lesssim\epsilon_w^2/w
\le Y\epsilon_w$.
Cauchy--Schwarz with the raw bound for $|X|^2$ bounds the mixed term
by $CY\epsilon_w$.  This proves \eqref{eq:D9}.

The largest $\epsilon_w$ is $E_0$, so summing dyads shows that
$\int_{x\le v\lesssim Y}D_v^2\,dv/Y\lesssim_\circ E_0$.
We shall also use the following first-moment consequence:
\begin{align*}
 Y\int_{v\asymp w,\ v\ge x}D_v\,dv
 &\lesssim_\circ Y(wY\epsilon_w)^{1/2}\\
 &=\frac{x^{3/2}Y^{7/4}w^{1/2}}
         {\lambda_1^{1/4}\max(w,U)^{3/4}}
 \lesssim_\circ x^3\lambda^{-C}Z^{7/4}.
\end{align*}
The expression in $w$ is largest at $w\asymp U$; substitution of
$Y=UZ$ and $U=x\lambda^{-s_*/2}$ proves the last inequality.

\end{proof}

\paragraph{Aliases and projected endpoints.}
We next specify how the estimates apply inside a larger finite topology. For a pair split at \(v\), its overlap may be used either at the terminal leaves or projected, i.e. using endpoint posterior means at later cuts where the branches are separate. These estimates also hold in a larger fixed topology. The error replacing any such overlap by \(|X_v|^2\), or \(|X_v|^2\) by \(2L_v\) along a single path through the fork, has squared \(L^2\) norm bounded by \(C\mathbf E[X^{\mathsf T} C_v X+\operatorname{tr} C_v^2]\), by branch independence and conditional contraction. For example, conditional independence at the split gives
\[
 \mathbf E[(Q-|X_v|^2)^2\mid\mathcal F_v]
       =2X_v^{\mathsf T}C_vX_v+\operatorname{tr}C_v^2,
 \qquad
 \mathbf E[(2L_v-|X_v|^2)^2\mid\mathcal F_v]
       =4X_v^{\mathsf T}C_vX_v.
\]
Projecting either endpoint to a later cut is conditional expectation
on an already separate branch, so conditional contraction only reduces
these bounds.  Thus \eqref{eq:D7}, \eqref{eq:D2} control these \emph{alias} errors, while \(D_v^2\) bounds errors relative to \(S_v\) without the replacements. Products assigned to one and the same fork have the \textbf{same} ancestor value \(|X_v|^2\).

\begin{lemma}[Centered polynomial integrals]\label{lem:statistics-centered}
Fix a dyad \(v\asymp r\), \(r\gtrsim Y\), and a deterministic polynomial \(P_v\) of total degree \(j\ge1\) in the centered inserted value. Assume its remaining deterministic powers are bounded by the corresponding powers of \(e(v)\). Then the centered integral using \(2L_v\) on one path satisfies
\begin{equation}
\left\|\frac1r\int [P_v(2L_v-S_v)-\mathbf EP_v(2L_v-S_v)]\,dv\right\|_2
 \lesssim_\diamond z_r^0 r Z_h
       \begin{cases}1& j=1,\\ m_r r^{j-2}& j\ge2.\end{cases}             \tag{D10}\label{eq:D10}
\end{equation}
Here the coefficient of degree \(i\) can be any deterministic function bounded by \(C e(v)^{j-i}\). \end{lemma}
\begin{proof}
Write $P_v(t)=\sum_{i=0}^j c_i(v)t^i$, with
$|c_i(v)|\le Ce(v)^{j-i}$.  On the event
$|2L_v-S_v|\le x^{-\epsilon'}e(v)$, differentiation gives
$|P'_v(2L_v-S_v)|\lesssim x^{-C_j\epsilon'}e(v)^{j-1}$.
Equations \eqref{eq:A2} and \eqref{eq:D7}, with $k=j-1$, therefore give
\[
 \left\|\frac1r\int(P_v-\mathbf EP_v)\,dv\right\|_2^2
 \lesssim_\diamond \frac1r\mathbf E\int W e^{2j-2}\,dv/r
 \lesssim_\diamond (z_r^0)^2Z_h^2r^2
 \begin{cases}1,&j=1,\\m_r^2r^{2j-4},&j\ge2.\end{cases}
\]
The factors of two in the argument of $P_v$ only change constants.
Taking square roots proves \eqref{eq:D10}.  For the truncation, first
fix the allowed loss $\epsilon'>0$, then a finite moment order large
enough that \eqref{eq:D2} makes the tail negligible; $W$ is bounded.
Choose the geometric slack in \eqref{eq:D2} last, small enough for that
moment order.  The same truncation of the other factors at a fork,
followed by Cauchy--Schwarz in $dv/r$, bounds its alias errors by the
same expression.

At \(x\le v\le C Y\) a centered \textbf{linear} integral of \(2L_v-S_v\) with bounded coefficients instead costs \(O_\circ(E_0)\) before any mass normalization, by \eqref{eq:D9} and \eqref{eq:A1}. When there are no other random factors in the expectation of that linear term, no alias error is required even with an initial density tilt: conditional averaging reduces the overlap to \(|X_v|^2\) and then \(2L_v\) exactly.

\end{proof}

\paragraph{Stopping, forests, and the order of truncation.}\label{par:statistics-tilt-transfer}
For a tilted active path on the stop, \eqref{eq:A4} represents
its expectation as an untilted expectation multiplied by a density
having every fixed moment required here.  H\"older therefore transfers
the integrated second moments and squared errors in
\eqref{eq:D7}--\eqref{eq:D10}.  The needed norm is slightly above
$L^2$.  Interpolation between its $L^2$ bound and its deterministic
bound loses an arbitrarily small power, since the target scale is
polynomial in $x$ and the path observables are bounded.
The same observation applies to the low linear centered integral.

In a finite product, first truncate the other factors using their
fixed moments from \eqref{eq:D2}; their deterministic envelope bounds
then cost only the chosen small power loss, and their discarded tails
are negligible.  These operations concern active subtrees or bounded
deterministic values at masked pairs, as specified in \eqref{eq:A4}.
They do not by themselves give a common marginal law when the topology
varies.  Such a law must first be constructed for the centered integral
and its retained tests.
For a forest, apply the change of measure separately on each active
component.  A bounded number of component densities has every required
fixed moment on the stop, so H\"older handles their product.  Neither
this transfer nor \eqref{eq:A2} is an assertion about a common untilted
root for the whole forest.  When other random path factors remain,
first realize the centered integral and those factors under one fixed
marginal law, apply the unconditional centered estimate, and only then
apply H\"older.  The existence of that marginal when a fork coordinate
varies is proved explicitly in Lemma~\ref{lem:statistics-low-sources}.

\begin{corollary}[Projection to a common low cut]\label{cor:statistics-low-projection}
One more consequence of \eqref{eq:A3}, \eqref{eq:D2} is the following: for \(Y\lesssim c'\) up to a sufficiently small fixed scale, a pair split well before \(c'\) (say mass \(\le c'/8\)) changes in overlap by at most \(O_{p,\circ}(z_{c'}^0 c')\) in every fixed moment upon projecting down from later cuts to mass \(c'\). \end{corollary}
\begin{proof}
Use future doubling layers up to a small macroscopic mass before the long terminal block, hitting the later layer in each product via \eqref{eq:A3} with a full gap after the fork (also for projected rather than terminal ends). This is the two-branch raw argument of \eqref{eq:D6}.
\end{proof}

\subsection{Coarse paths and multiplier changes}
\begin{lemma}[Coarse and switch contributions]\label{lem:statistics-coarse-switch}
For sufficiently small fixed $\kappa>0$, the direct-strength contribution
from $Y>x^\kappa$ satisfies the direct part of
Proposition~\ref{prop:increment}.  The two final-strength multiplier
changes also satisfy its required power-saving bound.
\end{lemma}
\begin{proof}
\leavevmode\par
\paragraph{The large-$Y$ event.}

Choose $\kappa>0$ small relative to the power gains below.
Markov's inequality and $\mathbf E_\gamma Z^4\lesssim_\circ1$ give
\[
 \mathbf P_\gamma(Y>x^\kappa)\lesssim_\circ U^4x^{-4\kappa}.
\]
We first bound the contribution at each parameter value on this event.
Here $U$ retains its cap $\min(1,x\lambda^{-s_*/2})$.  Since
$v\,dh/dv\gtrsim\lambda_1U^3$ for $v\ge x$, the source estimate
and $0\le S\le1$ give
\begin{equation*}
 \int_{v\ge x}\mathbf EW\,dv\lesssim e_1^2,
 \qquad e_1=\frac{x^3}{\sqrt{\lambda_1}U^{3/2}}.
\end{equation*}
With unit raw bounds, mass slides of width $e_1$ then cost
$O(e_1)$ in both their mean and centered parts.  The trace term on
$v\ge\max(e_1,x)$ costs $O(e_1)$ by division by $v$; below $e_1$
use the unit raw bound.  Cauchy--Schwarz on dyads handles the mixed
covariance term.  Hence
\[
 \int_{v\ge x}D_v^2\,dv\lesssim_\circ e_1.
\]

The dotted clocks in the strength derivative are $O_\circ(\lambda)$.
In the pressure difference from \eqref{eq:A1}, the portion $v<x$
therefore costs $O_\circ(\lambda x)$.  At the other splits expand
$Q^2$ around the common untilted mean $S_v$.  The quadratic centered
term uses the preceding $D_v^2$ estimate.  For the linear term,
conditional averaging replaces $Q-S_v$ first by $|X_v|^2-S_v$ and
then by $2L_v-S_v$ exactly; there is no alias error in its expectation.
The centered integral estimate from \eqref{eq:A1} costs $O_\circ(e_1)$.
Transfer these estimates to the stopped tilted law by \eqref{eq:A4}.
The stopped derivative is consequently
$O_\diamond(\lambda(x+e_1))$ at each parameter value.

After parameter averaging, apart from $x^{-4\kappa}$, the two costs
are bounded by
\[
 \lambda xU^4\le x^5\lambda^{1-2s_*},\qquad
 \lambda e_1U^4=x^3\lambda^{(1-\eta)/2}U^{5/2}
       \lesssim x^{3+(1-\eta)/s_*}.
\]
For the second bound, if $U=1$ then $\lambda\le x^{2/s_*}$ and the
claim follows directly.  If $U<1$, substitution gives
$x^{11/2}\lambda^{(1-\eta)/2-5s_*/4}$.
The exponent of $\lambda$ is negative for $s_*$ sufficiently close
to $1/2$; using $\lambda\ge x^{2/s_*}$ gives the same bound.
Since $s_*<1/2$ and $\eta$ is chosen small, both terms leave a fixed
power saving beyond $x^5$.  Choose $\kappa$ to preserve it after
$x^{-4\kappa}$, subpower losses, and strength integration.

\paragraph{The multiplier changes.}
In varying \(\chi_{\rm lin}\) and then \(\chi_e\) at final strength, \(U\) is power small. Use the direct active and plateau coordinates and tiles of \eqref{eq:D5} on macroscopic bands, also when \(Y>x^\kappa\). The deterministic width properties there only use \eqref{eq:D1}, \(S,q\) bounded, and do not require \(Y=o(1)\). The plateau lower slope in \eqref{eq:D3} comes from $\chi_eH_e$;
for $\chi_e>0$ its use therefore costs $\chi_e^{-1}$.  The active
part has at least the same weaker bound.  Consequently for \(H=H_{\rm lin}\) or \(H_e\), respectively,
\begin{equation*}
\mathbf E\int H^2 W\,dv\ \lesssim_\diamond\
   \chi_e^{-1}x^6\lambda_1^{-1} Z^2
       \begin{cases}
          x^6,\\ \lambda_1^2 U^4.
       \end{cases}
\end{equation*}
Indeed both are supported at masses bounded below, the first kernel is uniformly \(O(x^3)\), the second bounded by \(O_\diamond(\lambda_1\mathfrak h)\); apply the tile expectation bound of \eqref{eq:D7} with that same loss \(\chi_e^{-1}\) and the integrated squared width estimate. By \eqref{eq:A1}, \eqref{eq:A4} the derivative difference along each multiplier costs the square root of this bound on the stopped set (again exact linear conditional averaging); \(\chi_e^{-1/2}\) is integrable. Both \(x^6/\sqrt{\lambda_1}\) at final strength and \(x^3\sqrt{\lambda_1} U^2\) allow power saving beyond \(x^5\). This handles the multiplier changes.
\end{proof}


\subsection{Uniform coarse accuracy on size paths}
The pointwise trial excess in \eqref{eq:C2} provides an independent
starting estimate on size paths. We turn it into an integrated bound
for $|S-q|$ by a convex variation of the mass coordinate, then use the
scalar gap to make that bound uniform. The final moment upgrade uses
the same block and buffer identities as above with unit raw bounds.
Its purpose is to supply one positive exponent for every fixed moment
order before choosing a Taylor degree.

\Needspace{4\baselineskip}
\begin{proposition}[Uniform all-moment accuracy]\label{prop:statistics-size-accuracy}

Here \(\lambda=x^{p_*}\); assume \(p_*\) sufficiently small. Uniformly at every parameter value on smoothed size paths in the untilted system there is a fixed \(\delta_*>0\), available independently of further decreases in \(p_*\), such that
\begin{equation}
\sup_v\|Q-q(v)\|_p\le C_p x^{\delta_*}                  \tag{D11}\label{eq:D11}
\end{equation}
for a pair split at \(v\), for \textbf{every} fixed \(p\). \end{proposition}
\begin{proof}
Use both coordinates \(v,t=t(v)\). Up to \(t_M=M/\ell^2\), before the terminal constant piece beginning at \(v_M\), both coordinate derivatives are bounded by \(C/\omega\); the two are inverse, with \(v_M=\ell+O(\omega M)\). For \(v>x\) we have \(dh/dt\gtrsim\lambda v^2\) by \eqref{eq:P4}.

\paragraph{Freezing beyond the wall.}
At \(t\ge t_M/2\) and thereafter \(1-q\) and \(1-S\) are smaller than every power of \(x\). In fact for the scalar spin the conditional mean martingale above a fixed macroscopic time has diffusion coefficient between \(c(1-V_z^2)\) and \(C(1-V_z^2)\), giving exponential decay of \(\Gamma'\) as after \eqref{eq:P2}. Thus \(dh/dt\ge c>0\) above a sufficiently large fixed time (with \(b\) constant there). For each actual spin, the conditional mean \(m_i=\sqrt m (X_v)_i\) has martingale bracket at least \(c(1-m_i^2)^2dh\) there, by the \(i\)-th diagonal source-Hessian bound. Indeed the field contributes diffusion matrix \(m H_v\) to all these coordinates (the source-Hessian and clock identities in~\cite[Section~2]{PE}); matrix innovations are independent. More explicitly, since the masses here are bounded below,
$H_{ii}\ge c(1-m_i^2)/m$.  The field contribution to quadratic
variation is at least
$m^2H_{ii}^2dh\ge c(1-m_i^2)^2dh$; the independent matrix
innovations only increase it.  With $dh/dt\ge c$, localized It\^o
calculus for $F(u)=\sqrt{1-u^2}$ gives
\[
 \frac d{dt}\mathbf EF(m_i(t))
 \le -c\mathbf E
       \frac{(1-m_i(t)^2)^2}{(1-m_i(t)^2)^{3/2}}
 =-c\mathbf EF(m_i(t)).
\]
Thus its expectation decays exponentially, uniformly in $i,m$. The calculation works first across filled finite increments at masses bounded below, then by approximation. Averaging proves the assertion; since \(Q\le1,\ \nu Q=S\), it gives \eqref{eq:D11} in this region already.

\paragraph{Integrated fluctuation control.}
Let \(\mathfrak a=1/200,\ V_0'=x^{\mathfrak a}\), with \(p_*<\mathfrak a/100\). On \([V_0',v_M]\), \eqref{eq:A1} gives
\begin{equation*}
\int \mathbf E W\,dt\lesssim x^6/(\lambda (V_0')^3).
\end{equation*}
Let $R=x^6/(\omega\lambda(V_0')^3)$.  The bounds on the two
coordinate derivatives give
\[
 \int\mathbf EW\,dv
   =\int\mathbf EW\frac{dv}{dt}\,dt\lesssim R,
 \qquad
 \int\mathbf EW\Big(\frac{dt}{dv}\Big)^2dv
   =\int\mathbf EW\frac{dt}{dv}\,dt\lesssim R.
\]
The second integral is the one required by \eqref{eq:A1} for a
$t$-average, since that average has weight $dt/dv$ when expressed in
mass.  Mass slides and $t$-slides can therefore each use width
$R^{1/2}$, with unit raw bounds.  They give
\begin{equation}
\int_{V_0'}^{v_M} D_v^2 dv+
       \int_{t(V_0')}^{t_M} D_{v(t)}^2 dt\le x^{1/4}.               \tag{D12}\label{eq:D12}
\end{equation}
Use the bracket identities of Lemma~\ref{global:buffer-brackets}, with raw unit bounds, monotone \(S\), and width \((x^6/(\omega\lambda (V_0')^3))^{1/2}\) in each coordinate, using \(dt/dv\) in the \(t\)-averages in \eqref{eq:A1}. The \(\operatorname{tr} C^2\) terms cost at most \(C x^6/(\omega\lambda(V_0')^4)\) and Cauchy suffices for \(X^{\mathsf T}CX\) (coordinate lengths at most \(O_\circ(1)\)). All these bounds have the indicated slack.

\paragraph{A convex mass variation.}
We use the trial gap in \eqref{eq:C2} to control the integral
of $|S-q|$ in mass.  Let $\phi$ be a smooth bump of either sign,
supported between $V_0'$ and $\ell$, and set
\[
 \frac1{m_\theta(v)}=\frac1v-\theta\frac{\phi(v)}{v^2},
       \qquad 0\le\theta\le1.
\]
The bumps used below have small slope and $|\phi|/v\ll1$, making
$m_\theta$ an increasing bi-Lipschitz change of mass.  Reassign each
clock increment formerly at $v$ to $m_\theta(v)$, keeping the increment
itself, the terminal length, and the terminal self correction fixed.
Let $F(\theta)$ be the resulting pressure.

The exponential variational formula expresses each step power mean
as a conditional supremum of expectation minus relative entropy
divided by the step mass.  After iteration, the complete pressure is
a supremum over the choices of transition laws; the entropy
coefficients are the reciprocal masses $1/m_\theta(v)$, which are
affine in $\theta$.  Thus $F$ is convex.  The initial mass-zero
expectations are unchanged.  Passing from step clocks to the fixed-size
limit preserves convexity.  The segment formula \eqref{eq:A1} gives
\[
 F'(0)=\frac14\int(Bb'+2Sh')\phi\,dv.
\]
The clocks are Lipschitz on the bump at fixed size, so this chain rule
also follows directly by approximating with continuous observation cuts.

In the new mass coordinate $w=m_\theta(v)$, compare to the fixed scalar
trial $K(w)$.  Square completion along the affine interpolation to that
trial gives
\[
 F(\theta)\le f_{\rm spin}(K)+P(\theta),\qquad
 P(\theta)=\int\frac{(K(w)-h(v))^2}{4b(v)}\,dw,
       \quad v=m_\theta^{-1}(w).
\]
At $\theta=0$, \eqref{eq:C2} bounds this trial excess by $x^{4+c}$.
Convexity therefore yields
\[
 F'(0)\le F(1)-F(0)
       \le x^{4+c}+P(1)-P(0).
\]
Differentiating the penalty and returning to the old coordinate gives
\begin{equation*}
 P'(\theta)=\tfrac14\int(z_\theta^2b'+2z_\theta h')
                    \partial_\theta m_\theta(v)\,dv,
 \qquad z_\theta=\frac{K(m_\theta(v))-h(v)}{b(v)}.
\end{equation*}
If the oscillation of $K$ on the swept support is at most
$x^{3\mathfrak a}$, then $z_\theta-z_0=O(x^{3\mathfrak a})$.
Also
$\partial_\theta m_\theta=\phi(1-\theta\phi/v)^{-2}
=\phi[1+O(\sup|\phi/v|)]$.
These are the two errors in comparing $P'(\theta)$ to $P'(0)$.
Integrating that comparison in $\theta$ gives
\begin{equation}
\int [(B-z_0^2)b'+2(S-z_0)h']\phi
 \le 4x^{4+c}+C(x^{3\mathfrak a}+\sup|\phi/v|)\int(b'+h')|\phi|.       \tag{D13}\label{eq:D13}
\end{equation}
Here \(z_0=q-d/b\) is bounded and \(|d|\lesssim_\circ x^3\) throughout, by \eqref{eq:P6}.

Partition $[V_0',\ell]$ into cells of mass width
$l=x^{5\mathfrak a}$.  Give each cell a fixed enlargement, with
uniformly bounded overlap.  Discard the boundary cells and every cell
whose enlargement has $K$-oscillation greater than $x^{3\mathfrak a}$.
Since $K$ is nondecreasing and its total increase is $O(M)$,
\[
 \#\{\hbox{discarded interior cells}\}\,x^{3\mathfrak a}
       \lesssim M,
 \qquad
 |\hbox{discarded mass}|\lesssim l+Ml/x^{3\mathfrak a}.
\]
On each retained enlargement, Lipschitz continuity of $\Gamma$ gives
$q$-oscillation $O(x^{3\mathfrak a})=o(x^{\mathfrak a})$.
If $|S-q|\ge Cx^{\mathfrak a}$ somewhere in its cell, monotonicity
of $S$ preserves that sign, at half the threshold, on a one-sided
interval of length comparable to $l$ in the enlargement: use the
forward side for a positive deviation and the backward side for a
negative one.

Choose a smooth bump $\phi$ on this one-sided interval, with that sign,
height $cl$, slopes bounded by a sufficiently small fixed constant,
and $\int|\phi|\asymp l^2$.  Its swept support remains in the good
enlargement.  Moreover $|\phi|/v\lesssim l/V_0'=x^{4\mathfrak a}$;
the small slope and this bound make all $m_\theta$ increasing and
bi-Lipschitz.  Since $|d|\lesssim_\circ x^3$, $b\asymp1$, and
$q\gtrsim V_0'$ on this support, $z_0=q-d/b\ge0$ and
\[
 (S-z_0)\phi\gtrsim x^{\mathfrak a}|\phi|.
\]
Put $I_\phi=\int h'|\phi|$.  The error on the right of \eqref{eq:D13}
is at most
\[
 C(x^{3\mathfrak a}+x^{4\mathfrak a})(1+x^{-\mathfrak a})I_\phi
       =O(x^{2\mathfrak a})I_\phi
       =o(x^{\mathfrak a}I_\phi),
\]
because $b'\lesssim h'/V_0'$.

To estimate the remaining term on its left, use the exact identity
\[
 B-z_0^2=(S-z_0)(S+z_0)+D_v^2.
\]
The first summand has favorable sign after multiplication by
$b'\phi$.  The second can have unfavorable sign only for a negative
bump.  If $b'$ meets the bump, its support lies in the regular low
region, where $b'\lesssim\lambda$ and
$h'\gtrsim\lambda(V_0')^2$.  By \eqref{eq:D12},
\[
 \int D_v^2b'|\phi|\lesssim\lambda l x^{1/4},\qquad
 x^{\mathfrak a}I_\phi
       \gtrsim\lambda x^{3\mathfrak a}l^2.
\]
Their ratio is $O(x^{1/4-8\mathfrak a})=o(1)$.  If the bump lies
outside that low region, $b'=0$ and there is no variance error.
Finally $h'\gtrsim x^3$ everywhere on the bump: use the regular low
inverse before $4v_b$, and $H_{\rm lin}$ thereafter.  Thus
\[
 x^{\mathfrak a}I_\phi\gtrsim x^{3+11\mathfrak a}
       \gg x^{4+c},
\]
contradicting \eqref{eq:D13}.  Every retained cell therefore has
$|S-q|\lesssim x^{\mathfrak a}$.  Since $S,q$ are bounded, the low
interval of length $V_0'$ and the discarded cells contribute
$O(x^{\mathfrak a})$ and $O(l+Ml/x^{3\mathfrak a})$, respectively.
The strip between $\ell$ and $v_M$ has mass $O(\omega M)$, so its
contribution is $O(\omega M)=o(x^{\mathfrak a})$, including its part
before the frozen region.  The remaining tail is controlled by the
freezing estimate for $t\ge t_M/2$.  This proves
\begin{equation}
\int |S-q|\,dv\lesssim x^{\mathfrak a}.         \tag{D14}\label{eq:D14}
\end{equation}

\paragraph{From integrated means to pointwise means.}
Set \(K_S=h+b S\), nondecreasing, and write \(\Gamma_S\) for its scalar response (time variable), extending both profiles as needed by terminal annealed clocks. Use one common terminal extension for the two monotone clocks.
The area between their graphs can be integrated first in mass or first
in time, giving
\[
 \int|\alpha_{K_S}(t)-\alpha_K(t)|\,dt
       =\int|K_S(v)-K(v)|\,dv
       \lesssim x^{\mathfrak a}
\]
by \eqref{eq:D14} and $K_S-K=bf+d$.  Thus
\begin{equation*}
\sup_t|\Gamma_S(t)-\Gamma(t)|\lesssim (1+M)^C x^{\mathfrak a}
\end{equation*}
on the needed times. To see the polynomial constant, interpolate inverse profiles with evaluation time fixed; \(\Gamma(y)/2\) is the root impulse derivative in the profile. Its derivative uses the \hyperref[par:inverse-kernel-estimates]{post/post Hessian kernel} in the proof of Lemma~\ref{inv:sector-response}, which is uniformly \(O(1+M)\) for mixtures as well: the spatial derivative \(L_z^{\,y}(s)\) of the conditional scalar pair moment divided by two is bounded by the labeled source derivative rule (also directly by telescoping the conditional transition log-density when shifting the field, with bounded gradient of each scalar continuation). Impulse formulas pass from step profiles as there. Extend to one common terminal time \(O(M+1)\).

At reference $K_S$, the returning mismatch is exactly
$h+bQ-K_S=b(Q-S)$.  Apply \eqref{eq:A7} with $P=1,k=1$.
Its absolute attachment measure has only one prescribed fork: the
tested pair's split at $v$.  An insertion tied there costs $D_v$;
a freely placed insertion has bounded mass density and costs
$\int D_u\,du$.  The row-deletion contact contributes $O(m^{-1})$.
Consequently
\begin{equation*}
|S_v-\Gamma_S(K_S(v))|\lesssim D_v+\int D_u du+O(m^{-1}).
\end{equation*}
All row changes are bounded uniformly. Integrate over \(0<t<t_M\) at \(v=v(t)\), using \eqref{eq:D12}, discarding \(v<V_0'\) at cost \(O((1+M)V_0')\) since the low inverse is regular, and using the frozen estimate for post-wall attachments. Combine with the response comparison and \(|d|\lesssim_\circ x^3\); we get \(\int |S_{v(t)}-\Gamma(t+b(v)(S_v-q(v)))|dt\lesssim(1+M)^C x^{\mathfrak a}\). To use this residual estimate, write $f=S-q$ and integrate the
scalar derivative along the segment from $t$ to $t+bf$:
\[
 S-\Gamma(t+bf)
 =f\int_0^1[1-b\Gamma'(t+ybf)]\,dy.
\]
The average in brackets is at least $c\lambda v^2$ by
\eqref{eq:P5}.  On $v\ge x^{\mathfrak a/10}$, division by this gap
therefore gives, allowing the preceding subpower factors,
\begin{equation*}
 \int_{v(t)\ge x^{\mathfrak a/10}}|f(v(t))|\,dt
 \lesssim_\circ x^{4\mathfrak a/5-p_*}
 \le x^{\mathfrak a/2}.
\end{equation*}
Here $p_*<\mathfrak a/100$ leaves room for the weakened exponent.
The function $S_{v(t)}$ is nondecreasing, whereas $\Gamma(t)$ is
Lipschitz.  A deviation of height $d$ therefore persists at height at
least $d/2$ on a one-sided $t$-interval of length comparable to $d$.
The integral bound forces $d^2\lesssim x^{\mathfrak a/2}$, giving
$|f|\lesssim x^{\mathfrak a/4}$ above mass
$2x^{\mathfrak a/10}$.  Both-sided buffers fit until the frozen
region, which was treated separately.  Below this cutoff,
monotonicity and $q(v)\asymp v$ near zero give
$S_v+q(v)\lesssim x^{\mathfrak a/10}$.  Consequently
$\sup_v|S_v-q(v)|\lesssim x^{\mathfrak a/10}$.

\paragraph{The all-moment upgrade.}
Fix a finite moment order $p$.  First take a cut of mass
$v\ge4x^{\mathfrak a/10}$ and scalar time $t\le t_M/2$.
Its small $t$-neighborhoods stay above $V_0'$.  Apply Lemma~\ref{global:residual-blocks} with unit raw bounds.
For a later block at $t$-distance $l$, the gap has field length at
least $c\lambda(V_0')^2l$ and mass at least $V_0'$.  Its squared
projection coefficient and the integrated earlier anchor cost are
\[
 \frac{C_{p,\circ}x^6}{\lambda(V_0')^4l}
       \quad\hbox{and}\quad C_pl,
\]
respectively.  Stop the doubling layers at a fixed positive
$t$-distance and then use a terminal block.  Summing gives
\begin{equation*}
 \left\|\int W\,dt\right\|_{p/2}
       \lesssim_{p,\circ}\frac{x^6}{\lambda(V_0')^4}
\end{equation*}
on each bounded $t$-window under consideration.

Put $\rho=x^{\mathfrak a/4}$.  Choose a past buffer and a future
buffer of width comparable to $\rho$ on the two sides of a smaller
neighborhood of the cut.  They remain at mass at least $V_0'$ and
below $t_M$; any part past $t_M/2$ is already controlled by freezing.
The global mean bound and Lipschitz continuity of $\Gamma$ give,
for two buffer times $t_1,t_2$,
\[
 |S_{v(t_2)}-S_{v(t_1)}|
 \lesssim x^{\mathfrak a/10}+|t_2-t_1|
 \lesssim x^{\mathfrak a/10},
\]
since $|t_2-t_1|=O(\rho)$ and $\rho=x^{\mathfrak a/4}
=o(x^{\mathfrak a/10})$.  Applying \eqref{eq:A1} to a
normalized $t$-average, with $dt/dv\lesssim1/\omega$, gives
\[
 \|L_{\rm buffer}-\mathbf EL_{\rm buffer}\|_p
 \lesssim_{p,\circ}
       \frac{x^3}{\sqrt{\omega\lambda}(V_0')^2\rho}.
\]
This is smaller than $x^{\mathfrak a/10}$ for our parameter choices.
The two-sided brackets and the common-past estimate in
Lemma~\ref{global:buffer-brackets} consequently give
\[
 \||X_t|^2-q(v(t))\|_p
       +\||X_t-X_{t'}|^2\|_p
       \lesssim_{p,\circ}x^{\mathfrak a/10}
\]
for cuts $t,t'$ in the smaller neighborhood.

For an overlap split there, project both endpoints to a cut a fixed
fraction of $\rho$ later.  The two projected branches have the preceding
norm and increment bounds, so polarization about their common ancestor
controls their inner product.  Equation \eqref{eq:A3}, now even with
unit anchors, bounds the residual projection errors using that
$\rho$-length gap.  For a smaller split, choose instead a central cut
at mass a sufficiently large fixed multiple of $x^{\mathfrak a/10}$.
Its lower comparable mass interval lies strictly after the fork and
supplies a full projection gap; it has the buffers just constructed.
Its squared posterior norm has $L^p$ size $O_p(x^{\mathfrak a/10})$.
Together with the frozen region, these estimates prove \eqref{eq:D11},
for example with $\delta_*=\mathfrak a/50$.  The stopped transfer
\eqref{eq:A4} gives the same coarse accuracy for shifted active overlaps.
\end{proof}

\begin{corollary}[A deterministic cap on the low scale]\label{cor:statistics-cap}
For \eqref{eq:D1} and the low statistics on size paths we may cap
$Y$ at $x^{\delta_*/3}$, decreasing $p_*$ if necessary, and still write
$Y=UZ$.  Indeed \eqref{eq:D11} gives
$|S_v-q(v)|^2\lesssim x^{2\delta_*}$, whereas on every bounded band
\[
 x^{2\delta_*}\lesssim x^{4\delta_*/3}/r^2.
\]
Thus the capped scale still satisfies \eqref{eq:D1}.  Its new majorant
$Z=\min(Z_{\rm old},x^{\delta_*/3}/U)$ is at least one and no larger
than the old majorant, so its fourth parameter moment is preserved.


\end{corollary}

\subsection{High-mass closure at final strength}
The coarse estimate fixes the Taylor order in the cavity equation.
The sharper low-mass integral bounds then control its sources below a
cutoff, while the pair inverse controls the stochastic linear terms
above that cutoff. Keeping one prescribed high pair as a centered test
will preserve its second norm through this calculation.

\begin{proposition}[High-mass means and variances]\label{prop:statistics-high}

Take \(Z\) as at the end of \eqref{eq:D11}; by \eqref{eq:C2} take also \(Z_g\ge1,\ \mathbf E_\gamma Z_g^2\lesssim1\) with \(\int v^2\nu[(Q-q)^2]\,dv\le x^{4+c} Z_g^2\) untilted. Choose
\begin{equation*}
V_i=x^{a_i}\quad(0\le i\le2),\qquad
 p_*\ll a_2\ll a_1\ll a_0\ll\min(c,\delta_*).
\end{equation*}
There are fixed \(b_H>0,C_H<\infty\), with \(b_H\gg C_Ha_1,\ a_1\gg C_Ha_2\), for which
\begin{equation}
\sup_{v\ge V_1}|f(v)|\le x^{2+b_H}(Z^2+Z_g),\qquad
 D_{H,i}:=\sup_{v\ge V_i}D_v\le x^{3-C_H a_i}Z^2\quad(i=1,2).       \tag{D15}\label{eq:D15}
\end{equation}

\end{proposition}

We first isolate the centered low-source estimate needed to invert the
high equations.  Its topology assumptions are part of the statement.
\begin{lemma}[Low sources with a retained high pair]\label{lem:statistics-low-sources}
Consider an allocation with bounded degree and bounded topology
weights.  The old prescribed topology is either absent or consists of
one pair splitting at mass at least $V_i$; in particular it contains no
prescribed low fork.  Call designated factors at $v<V_i$ low.
Each low factor is either stochastic, of the form $\delta=Q-S_v$,
or deterministic, bounded by $O_\diamond(e(v))$, with $e=Y$ for
$v\lesssim Y$.  There is at least one stochastic low factor and no
stochastic factor at a higher mass.  Deterministic factors at higher
masses may be included in the coefficient.  If there is exactly one
low factor, assume that its coefficient is $O_\diamond(v)$.
Projected pairs whose endpoints are already on separate branches are
allowed, with the convention following \eqref{eq:D9}.

There are two conclusions.  In the untilted law, test the product by
$P=Q_{ab}-S_{ab}$ for the prescribed high pair.  Alternatively, compare
its tilted and untilted expectations on the stop, with no $P$, using
$\delta^*=Q^*-S_v$ in the tilted law and the same deterministic
coefficients in both laws.  After signed integration the respective
bounds are
\begin{equation}
 O_\diamond(x^3\lambda^{-C}Z^2){\cal P}+O(x^K),
 \qquad {\cal P}=D_{ab}\ \hbox{or respectively }1.                 \tag{D16}\label{eq:D16}
\end{equation}
The negligible error exponent $K$ can be arbitrarily large.

Extra auxiliary slots are permitted.  On each fixed tree shape the
absolute joint density of freely placed distinct fork vertices is
bounded; deterministic weights and restrictions may depend on all
their times.  An additional multiplier depending only on the high
times can be factored out and integrated separately.  Ties to a low
vertex reuse its freely integrated coordinate, by the signed rule;
no initially free low fork requires an atomic time measure.
This distinction concerns the underlying allocation measure.  The old
high fork may carry its prescribed atom.  A new low fork is instead a
free coordinate with bounded density; later attachments to it reuse that
coordinate.  Merely having the same mass as a different vertex supplies
no additional atom.  The argument below retains this distinction before
discarding deterministic restrictions in absolute estimates.

\end{lemma}
\begin{proof}
\leavevmode\par
\paragraph{The region below $Y$.}
At $v<x$, a stochastic insertion costs $O_p(Y)$.  If it is
the only low factor, the required coefficient is $O(v)$, and the
integral is $O(x^2Y)$.  A second factor at the same vertex gives
$O(xY^2)$ instead.  At a distinct low vertex its integrated envelope
is at most $O_\diamond(Y^2)$, giving $O_\diamond(xY^3)$.
All three are bounded by $O_\diamond(x^3\lambda^{-C}Z^2)$, since
$Y=x\lambda^{-s_*/2}Z\le1$.  Additional factors are bounded after
truncation.

For $x\le v\lesssim Y$, retain one stochastic factor in second norm.
In the sole-factor case its coefficient is at most $CY$; in the tied
case the other envelope is $O(Y)$; in the distinct-vertex case its
integrated envelope is $O_\diamond(Y^2)\le O_\diamond(Y)$.
All three cases are therefore bounded by
\[
 Y\int_{x\le v\lesssim Y}D_v\,dv
       \lesssim_\circ x^3\lambda^{-C}Z^{7/4}
       \le x^3\lambda^{-C}Z^2,
\]
using \eqref{eq:D9} and summing dyads.  Cauchy--Schwarz with $P$
supplies its factor $D_{ab}$; the remaining factors are truncated at
their deterministic envelopes.  Tilted active moments transfer on the
stop as already specified.

\paragraph{Alias errors above $Y$.}
For stochastic insertions at $v\gtrsim Y$, first replace $Q-S_v$ by
$|X_v|^2-S_v$.  On a dyad $r$, the second line of \eqref{eq:D7}, with
$Z_h=1$, gives the following integral of the alias's second norm:
\[
 A_r:=\int_{v\asymp r}\|Q-|X_v|^2\|_2\,dv
       \lesssim_\diamond z_r^0r^2
       =\frac{x^3}{\sqrt{\lambda_1}r}.
\]
The sole-factor coefficient cancels this denominator: $rA_r\lesssim
x^3/\sqrt{\lambda_1}$.  If a second factor is tied to the same fork,
Cauchy--Schwarz in mass uses its integrated envelope budget and gives
\[
 \int_{v\asymp r}e(v)\|Q-|X_v|^2\|_2\,dv
       \lesssim_\diamond z_r^0r^2m_r
       =\frac{x^3Y^2}{\sqrt{\lambda_1}r^2}
       \le\frac{x^3}{\sqrt{\lambda_1}}.
\]
For a second distinct vertex, integrate its envelope separately; the
cost is at most $A_r\int_{v<V_i}e(v)dv
\lesssim_\diamond x^3Y^2/(\sqrt{\lambda_1}r)$, which is sufficient
because $r\gtrsim Y$ and $Y\le1$.  Truncation handles any further
factors, and Cauchy--Schwarz with $P$ again retains $D_{ab}$.

\paragraph{Centered integration in a fixed marginal.}
Process the remaining stochastic vertices in decreasing time
order.  At the current vertex, replace every tied overlap by the same
ancestor value $|X_v|^2-S_v$, and then replace its power by the
corresponding power of $2L_v-S_v$ on one continuation through that
vertex.  The alias estimates just proved bound these changes.  Next
replace this polynomial by its deterministic \emph{untilted} mean.
The mean is bounded by the corresponding envelope power.

The centered error in this last replacement is estimated by
\eqref{eq:D10}.  Its bound has the same three forms as the alias
bounds: a sole linear factor uses the extra coefficient $v$; a tied
product uses the $m_r$ gain; a factor at a distinct vertex uses the
integrated envelope there.  Earlier random ancestor factors are
truncated at their deterministic envelopes before this estimate.
To apply \eqref{eq:D10}, we must realize the entire moving-$v$
integral and its retained random tests under a single fixed marginal
law.  The next construction verifies that requirement.

More explicitly, fix the tree shape, its time order, every other distinct
fork coordinate, and every projection cut.  If an old leaf descends
through the moving vertex, retain that leaf as the continuation defining
$L_v$.  Both leaves of the old pair then lie in the same child until
their prescribed higher split.  The retained old pair consequently has
its original, fixed split law.  If no old leaf descends through the
moving vertex, the chosen continuation separates from the retained old
tree at a different, already fixed vertex.  In both cases normalized
continuations of all other children integrate to one.  Earlier random
ancestor factors are evaluated before the moving coordinate; later
stochastic factors have already been replaced by deterministic means.
Projected old endpoints are evaluations of the same retained paths at
fixed later cuts.  Auxiliary labels with no remaining spin observable
can be deleted by the same consistency rule.

Thus the entire varying-$v$ integral lives under one marginal law.
We do not condition the centered estimate on the retained random test
$P$.  Let $I$ denote the centered integral, and truncate the product $R$
of the earlier factors at its deterministic envelope product $M_R$.
Then
\[
 |\mathbf E[PRI]|\le M_R\|P\|_2\|I\|_2
       =M_R D_{ab}\|I\|_2.
\]
Equation \eqref{eq:D10} is applied to $I$ under the fixed marginal,
before this Cauchy--Schwarz inequality.  The discarded tails are
negligible by the previously specified moment truncation.  In the
tilted comparison, transfer and H\"older require a norm only slightly
above $L^2$, with the already permitted small power loss.  A stochastic tie at the moving vertex
is one polynomial power of its ancestor value; it does not add an
independent coordinate or create a second atom. In the comparison version all moving stochastic times here are active; the same fixed marginal reasoning and \eqref{eq:A4} suffice. After successive replacements (split ties as one power), the deterministic product vanishes against centered \(P\), or cancels across ensembles. This proves \eqref{eq:D16}.
\end{proof}

\begin{remark}[An earlier pair retained above a moving fork]\label{rem:statistics-earlier-pair}
The same marginal construction will be needed in Section~\ref{sec:expansion}
when a retained pair splits at a fixed mass $u<v$, rather than above the
moving fork.  At most one endpoint of that pair can descend through $v$.
Choose its path for $L_v$, if such an endpoint exists, and retain the
other endpoint.  After the alias replacement the other child at $v$ is
unneeded and can be pruned.  The retained paths still split at $u$, so
their law is independent of $v$.  Even if the earlier pair has already
been projected to a fixed cut $c>v$, its projected endpoints are
fixed-time observables of these same retained paths.  The centered
inequality remains unconditional, followed by H\"older with the earlier
random factor.  Figure~\ref{fig:fixed-marginal} records this configuration.
\end{remark}

\begin{figure}[htbp]
\centering
\begin{tikzpicture}[>=Latex,x=1cm,y=1cm,every node/.style={font=\small},dot/.style={circle,fill,inner sep=1.6pt}]
\node[dot,label=above:$u$] (u) at (0,2.6) {};
\node[dot,label=left:$v$] (v) at (-.9,1.3) {};
\coordinate (a) at (-1.7,0);
\coordinate (b) at (1.3,0);
\coordinate (q) at (-.1,0);
\draw[thick] (u)--(v)--(a) node[below] {$a$};
\draw[thick] (u)--(b) node[below] {$b$};
\draw[dashed] (v)--(q) node[below,align=center] {extra\\branch};
\node at (0,-1.15) {Before marginalization};
\draw[->] (2.1,1.4)--(3.7,1.4) node[midway,above,align=center,font=\footnotesize] {replace aliases;\\prune};
\node[dot,label=above:$u$] (u2) at (5.5,2.6) {};
\coordinate (a2) at (3.8,0);
\coordinate (b2) at (6.8,0);
\draw[thick] (u2)--(a2) node[below] {$a$};
\draw[thick] (u2)--(b2) node[below] {$b$};
\node[draw,circle,inner sep=2pt,label=left:$v$] at ($(u2)!.5!(a2)$) {};
\node at (5.5,-1.15) {The retained law has the fixed split $u$};
\end{tikzpicture}
\caption{The earlier-pair configuration in the centered integration argument. The pair $(a,b)$ splits at the fixed mass $u<v$. Replacing the aliases and integrating out the extra branch leaves the two paths on the right. Their joint law is independent of $v$, which now marks an evaluation point rather than a retained fork. Values at fixed projection cuts remain observables of these paths. The single-path estimate is applied under this fixed law before H\"older's inequality handles the other random factors.}
\label{fig:fixed-marginal}
\end{figure}


\Needspace{6\baselineskip}
\begin{proof}[Proof of Proposition~\ref{prop:statistics-high}]
\leavevmode\par
\paragraph{The covariance equation.}
Fix an old pair $ab$ splitting at or above $V_i$, where
$i\in\{1,2\}$, and put $P=Q_{ab}-S_{ab}$.  For every tested pair
$jk$ at or above the same threshold, define
\[
 u_{jk}=\nu[P(Q_{jk}-S_{jk})].
\]
These are pair functions on extensions of the old labels $ab$; unused
labels are ignored by marginalization.  Apply \eqref{eq:A7} at the
reference $K$, expanding $J=b\delta+bf+d$, where
$\delta=Q-S$.  Choose a fixed truncation degree large enough that
\eqref{eq:D11} makes the true Taylor remainder $O(x^K)$ for any needed
fixed $K$.  The row-deletion error remains $O(m^{-1})=O(x^6)$.
Every all-deterministic term vanishes against $P$.

Retain the high stochastic linear term on the left.  The equations
have the form
\[
 u_{jk}-\sum_{cd:\ v_{cd}\ge V_i}{\bf M}_{jk,cd}b_{cd}u_{cd}
       =R_{jk}^{\rm low}+R_{jk}^{\rm nonlinear}
                         +O(x^K)+O(x^6).
\]
The first source includes the stochastic linear term below $V_i$;
the second contains at least two insertions.  This left side is exactly
\eqref{eq:H1}, because the positive-law expectation of each inserted
pair with $P$ depends only on the indicated labels.  By \eqref{eq:P5},
its gap on retained levels is at least $c\lambda V_i^2$.  Proposition
\ref{inv:pair-inverse} therefore inverts it with cost $O(x^{-Ca_i})$,
once $p_*\ll a_i$ is chosen.  Evaluate the result at $jk=ab$.
The inverse adds only a bounded number of freely attached labels.
Their attachment densities and existing-fork weights are bounded after
removing the inverse cost; they do not turn a new low fork into a
prescribed atom.

As explained in the \hyperref[par:inverse-limits]{fixed-resolution limit argument}, these operations can first be made on resolutions keeping the tested cut and threshold. At fixed system size gaps persist on fine approximations; inverse domination with block restrictions allows replacing path expectations and coefficients outside that inverse by continuous-path ones with errors tending to zero (empty step intervals can be refined for attachment). Our bounds apply to bounded measurable topology weights from that inverse after its cost factor is removed.

Consider a nonlinear term containing a high stochastic
factor $\delta_{cd}$.  Keep $P$ and that factor untruncated.  There is
at least one further small factor: truncate it at $x^{\delta_*/2}$,
and bound all other factors uniformly.  Equation \eqref{eq:D11}
controls the discarded tails in arbitrarily high fixed moments.
Cauchy--Schwarz on the two untruncated factors gives
\[
 |\nu[P\delta_{cd}\,\hbox{remaining factors}]|
 \lesssim x^{\delta_*/2}\|P\|_2\|\delta_{cd}\|_2+O(x^K)
 \le x^{\delta_*/2}D_{ab}D_{H,i}+O(x^K).
\]
This preserves the exact second-moment factors; no higher norm
proportional to $D_{ab}$ is assumed.  Every remaining source contains
only low stochastic factors, so Lemma~\ref{lem:statistics-low-sources}
applies, including to the low linear term.

\paragraph{The coefficient of a sole low factor.}
If just one low insertion occurs at \(v\), its scalar coefficient indeed costs \(O(v)\): in every term of the finite quotient it belongs to a spin pair moment (possibly with other pairs), and all other pairs in that moment, including the tested pair if present, are high. Immediately across the low branching its two ends thus give two different groups of odd counts. Conditional on prefixes up to that cut, group expectations on independent continuations are odd smooth bounded functions (for the odd groups) with bounded gradients by the scalar source rules. Their product costs \(O(K(v))=O(v)\) using Brownian moments and bounded drift. More explicitly, the two odd group functions $F_1(z),F_2(z)$
vanish at zero and have uniformly bounded derivatives, so
$|F_1(z)F_2(z)|\le C|z|^2$.  The shared scalar prefix field at time
$K(v)$ has second moment $O(K(v))$ by its Brownian representation
with bounded drift.  Hence this scalar coefficient is $O(K(v))=O(v)$.
Additional labels and coincidences within the groups do not change
the odd parity or the derivative bound.

Combining these estimates after inversion gives
\begin{equation*}
 D_{ab}^2\lesssim_\diamond x^{-Ca_i}
       \bigl(x^{\delta_*/2}D_{ab}D_{H,i}
              +x^3\lambda^{-C}Z^2D_{ab}+x^6\bigr).
\end{equation*}
The final $x^6$ is the row-deletion contact error.  Increasing the
Taylor degree does not remove it.  Take the supremum over $ab$ and
write $D=D_{H,i}$.  Choose the power losses so that
$Cx^{\delta_*/2-Ca_i}<1/2$.  Absorption then leaves
\[
 D^2\lesssim_\diamond x^{3-Ca_i}\lambda^{-C}Z^2D
                         +x^{6-Ca_i}.
\]
If $D^2\le AD+B$ with $D\ge0$, then $D\le A+\sqrt B$ up to a fixed
constant.  Since $Z\ge1$ and $p_*\ll a_i$, the factors
$\lambda^{-C}$ and the permitted small losses can be included in a
fixed $C_Ha_i$.  This proves
$D_{H,i}\le x^{3-C_Ha_i}Z^2$ for $i=1,2$.

\paragraph{The longitudinal equation.}
For $r\ge V_1$, apply \eqref{eq:A7} without a random test.
The longitudinal identity \eqref{eq:H2} separates the linear term
into its diagonal and its integral part.  The resulting scalar equation
is
\[
 [1-b(r)\Gamma'(K(r))]f(r)
     =\Gamma'(K(r))d(r)+R_{\rm int}(r)+R_{\rm nonlinear}(r)+O(x^6).
\]
The diagonal is in field units; returning from normalized mass changes
only bounded constants.  Its lower bound is
\[
 1-b(r)\Gamma'(K(r))\ge c\lambda V_1^2.
\]
For the integral term, the post kernel bound in Section~\ref{sec:inverse}
gains the earlier time $K(u)\lesssim u$ when $u$ is small.  It gives
\[
 |R_{\rm int}(r)|\lesssim(1+M)^C\left(
       \int|d|\,du+\int_{u<V_0}u|f(u)|\,du
                          +\int_{u\ge V_0}|f(u)|\,du\right).
\]
The low dyadic envelope budget and the raw bound below $Y$ imply
$\int_{u<V_0}u|f(u)|du\lesssim_\diamond V_0Y^2$.
For the last integral use \eqref{eq:C2} and the majorant $Z_g$:
Cauchy--Schwarz, followed by a harmless weaker cutoff bound, gives
$\int_{u\ge V_0}|f(u)|du\lesssim
x^{2+c/2-a_0}Z_g$.  Thus the three terms on this right side are
\[
 O_\circ(x^3),\qquad
 O_\diamond(x^{2+a_0-p_*s_*}Z^2),\qquad
 O(x^{2+c/2-a_0}Z_g).
\]

For a nonlinear allocation with a factor at the prescribed fork $r$,
keep that factor and truncate one of the others at
$x^{\delta_*/2}$.  Its contribution is
\[
 O(x^{\delta_*/2})\bigl(|f(r)|+x^3+D_{H,1}\bigr).
\]
If no factor is pinned at $r$, integrate one factor in second norm.
The low envelope and the high comparison budget give
\[
 \int\|Q-q(u)\|_2\,du+\int|d(u)|\,du
 \lesssim_\diamond Y^2+x^{2+c/2-a_0}Z_g+x^3.
\]
A further factor contributes $x^{\delta_*/2}$ after truncation, even
when it is tied to the same vertex.  The negligible tails use a fixed
higher moment from \eqref{eq:D11}.

Divide by $c\lambda V_1^2$.  The coefficient of $|f(r)|$ is at most
$Cx^{\delta_*/2-p_*-2a_1}$, up to the fixed polynomial and small
power losses, and is absorbed.  The remaining leading exponents before
division are
\[
 3,\qquad 2+a_0-p_*s_*,\qquad
 2+c/2-a_0,\qquad 2+\delta_*/2-p_*s_*.
\]
The already proved variance bound controls the extra
$x^{\delta_*/2}D_{H,1}$ term.  Choose $a_0$ small relative to
$c,\delta_*$, then choose $b_H$ to be a sufficiently small fixed
fraction of $a_0$.  Next take $a_1$ sufficiently small, choose
$a_2\ll a_1$, and only then take $p_*$ sufficiently small that the
loss from the gap, $(1+M)^C$, and every permitted small power loss is
less than the remaining margins above $2+b_H$.
This also ensures $b_H\gg C_Ha_1$.  We obtain uniformly for $r\ge V_1$
\[
 |f(r)|\le x^{2+b_H}(Z^2+Z_g),
\]
as required.  This choice of $a_2\ll a_1$, followed by $p_*$,
preserves both this bound and the covariance absorption.
\end{proof}

\begin{remark}[The powers of the parameter majorants]\label{rem:statistics-majorants}
The parameter moment used throughout is $\mathbf E_\gamma Z^4\lesssim_\circ1$.
The direct construction uses $e=\min(r,Z^2\mathfrak h)$, whereas its
low-region size analogue incorporates $Y$ in the width and uses
$e=\min(r,\mathfrak h)$.  This is the reason for $Z_h=Z$ in the direct
integrated bounds and $Z_h=1$ in the low-region bounds: no extra power
of $Z$ is implicit in $\lesssim_\diamond$.  In particular the right side
of \eqref{eq:D8} is $x^{-C_0\delta'}z_1^0 U^2Z^4$ up to the permitted
losses.  The lower-band factor $Z^{7/4}$ is bounded by $Z^2$, and the
low-source and high-variance estimates require only $Z^2$.
The high-mean estimate additionally uses $Z_g$ with
$\mathbf E_\gamma Z_g^2\lesssim1$.  Hence the combinations needed later,
$Z^4$ and $Z_gZ^2$, are integrable by Cauchy--Schwarz.  These powers are
also valid for path parameters and site counts that vary measurably:
the majorants and all windows are chosen pointwise from deterministic
path data, before taking a positive-law expectation.
\end{remark}

\paragraph{Parameter order.}
The final-strength saving $c$ from \eqref{eq:C2} and the exponent $\delta_*$ in
\eqref{eq:D11} are fixed before the final choice of $p_*$.  Hence the row Taylor
degree in the covariance equation is fixed using $\delta_*$, and can
make the remainder, apart from $m^{-1}$ contacts, smaller than $x^6$ by
an arbitrary fixed power.  The inverse then has a fixed polynomial loss
in $(\lambda V_i^2)^{-1}$ and $1+M$.  Choose $a_0$ relative to
$c,\delta_*$, next $a_1$ relative to $a_0$, next $a_2$ relative to $a_1$,
and only then $p_*$ relative to all these gaps.  This gives
$b_H\gg C_Ha_1$ and $a_1\gg C_Ha_2$ without a degree depending on a later
choice of $p_*$.  The local strengthening depths and their power slacks
may then be fixed as in Section~\ref{sec:parameters}.  The geometric
loss exponent used in the width construction is chosen last, after the
finite moment and truncation orders have been specified.
